% =========================================================
%  Foundations of Machine Learning
%  Lecture Notes by Dat Quoc Truong
%  A4 · pdfLaTeX · Overleaf-compatible
% =========================================================
\documentclass[11pt,a4paper,oneside]{book}

% ── Core packages ─────────────────────────────────────────
\usepackage[T1]{fontenc}
\usepackage[utf8]{inputenc}
\usepackage{palatino}
\usepackage{mathpazo}
\usepackage[a4paper,
            top=2.8cm, bottom=2.8cm,
            left=3.0cm, right=2.8cm,
            headheight=15pt]{geometry}

% ── Math ──────────────────────────────────────────────────
\usepackage{amsmath,amssymb,amsthm,mathtools,bm}

% ── Graphics & colour ─────────────────────────────────────
\usepackage{graphicx}
\usepackage[dvipsnames,svgnames,x11names]{xcolor}
\usepackage{pgfplots}
\pgfplotsset{compat=1.18}
\usepackage{tikz}
\usetikzlibrary{
  arrows.meta, positioning, shapes.geometric,
  calc, decorations.pathreplacing, fit, patterns,
  backgrounds, shadows}

% ── Tables ────────────────────────────────────────────────
\usepackage{booktabs,array,tabularx,longtable,multirow}
\renewcommand{\arraystretch}{1.25}

% ── Captions & floats ─────────────────────────────────────
\usepackage{caption,subcaption,float}
\captionsetup{
  font=small, labelfont=bf,
  skip=6pt, margin=10pt}

% ── Code listings ─────────────────────────────────────────
\usepackage{listings}
\definecolor{lst@bg}{RGB}{250,250,252}
\definecolor{lst@kw}{RGB}{0,64,153}
\definecolor{lst@str}{RGB}{163,21,21}
\definecolor{lst@cmt}{RGB}{100,100,100}
\lstset{
  language=Python,
  backgroundcolor=\color{lst@bg},
  basicstyle=\ttfamily\small,
  keywordstyle=\color{lst@kw}\bfseries,
  stringstyle=\color{lst@str},
  commentstyle=\color{lst@cmt}\itshape,
  numberstyle=\tiny\color{gray},
  numbers=left, numbersep=8pt,
  breaklines=true,
  frame=single, framerule=0.4pt,
  rulecolor=\color{gray!40},
  showstringspaces=false,
  tabsize=4, captionpos=b,
  xleftmargin=1.6em, framexleftmargin=1.6em,
  aboveskip=8pt, belowskip=6pt,
  morekeywords={Sequential,Dense,model,compile,fit,predict,
                Adam,BinaryCrossentropy,SparseCategoricalCrossentropy,
                MeanSquaredError,from_logits,activation,units,
                optimizer,loss,epochs,learning_rate,
                tf,keras,tensorflow,StandardScaler,
                XGBClassifier,XGBRegressor,np,plt}}

% ── Algorithms ────────────────────────────────────────────
\usepackage{algorithm,algpseudocode}

% ── Typography & spacing ──────────────────────────────────
\usepackage[final]{microtype}
\usepackage{setspace}
\setstretch{1.15}
\usepackage{parskip}
\setlength{\parskip}{5pt plus 1pt minus 1pt}
\setlength{\parindent}{0pt}
\usepackage{enumitem}
\setlist{itemsep=2pt,topsep=3pt}

% ── Headers & footers ─────────────────────────────────────
\usepackage{fancyhdr}
\pagestyle{fancy}
\fancyhf{}
\fancyhead[L]{\small\itshape\leftmark}
\fancyhead[R]{\small\thepage}
\renewcommand{\headrulewidth}{0.4pt}
\fancypagestyle{plain}{%
  \fancyhf{}
  \fancyhead[R]{\small\thepage}
  \renewcommand{\headrulewidth}{0pt}}

% ── Section formatting ────────────────────────────────────
\usepackage{titlesec}
\titleformat{\chapter}[display]
  {\normalfont\LARGE\bfseries}
  {\chaptertitlename\ \thechapter}{14pt}{\Huge}
\titlespacing*{\chapter}{0pt}{20pt}{18pt}
\titleformat{\section}{\large\bfseries}{\thesection}{1em}{}
\titlespacing*{\section}{0pt}{14pt}{5pt}
\titleformat{\subsection}{\normalsize\bfseries}{\thesubsection}{1em}{}
\titlespacing*{\subsection}{0pt}{10pt}{4pt}
\titleformat{\subsubsection}{\normalsize\itshape\bfseries}{\thesubsubsection}{1em}{}
\titlespacing*{\subsubsection}{0pt}{8pt}{3pt}

% ── Table of contents ─────────────────────────────────────
\usepackage{tocloft}
\setcounter{tocdepth}{2}
\renewcommand{\cftchapfont}{\bfseries}
\renewcommand{\cftchapleader}{\cftdotfill{\cftdotsep}}
\setlength{\cftbeforechapskip}{5pt}

% ── Theorem environments ──────────────────────────────────
\theoremstyle{definition}
\newtheorem{definition}{Definition}[section]
\newtheorem{example}{Example}[section]
\newtheorem{remark}{Remark}[section]
\theoremstyle{plain}
\newtheorem{theorem}{Theorem}[section]
\newtheorem{proposition}{Proposition}[section]

% ── Hyperref (load last) ──────────────────────────────────
\usepackage[
  colorlinks=true,
  linkcolor=black,
  urlcolor=NavyBlue,
  citecolor=black,
  pdftitle={Foundations of Machine Learning},
  pdfauthor={Dat Quoc Truong}]{hyperref}

% ── Colour palette for figures ────────────────────────────
\definecolor{pBlue}{RGB}{31,119,180}
\definecolor{pOrange}{RGB}{214,95,14}
\definecolor{pGreen}{RGB}{44,160,44}
\definecolor{pRed}{RGB}{192,0,0}
\definecolor{pGray}{RGB}{130,130,130}
\definecolor{pPurple}{RGB}{148,103,189}
\definecolor{pBrown}{RGB}{140,86,75}
\definecolor{pPink}{RGB}{227,119,194}
\definecolor{cA}{RGB}{70,130,180}
\definecolor{cB}{RGB}{210,80,80}
\definecolor{cC}{RGB}{60,165,90}
\definecolor{cD}{RGB}{220,140,0}
\definecolor{nodeBlue}{RGB}{173,216,230}
\definecolor{nodeGreen}{RGB}{198,230,198}
\definecolor{nodeGray}{RGB}{220,220,220}

% ── Math convenience macros ───────────────────────────────
\newcommand{\R}{\mathbb{R}}
\newcommand{\vx}{\vec{x}}
\newcommand{\vw}{\vec{w}}
\newcommand{\va}{\vec{a}}
\newcommand{\pd}[2]{\dfrac{\partial #1}{\partial #2}}
\newcommand{\IG}{\mathrm{IG}}
\DeclareMathOperator*{\argmin}{arg\,min}
\DeclareMathOperator*{\argmax}{arg\,max}

% =========================================================
\begin{document}
% =========================================================

% ── Title page ────────────────────────────────────────────
\begin{titlepage}
\centering
\vspace*{2.5cm}
{\fontsize{40}{48}\selectfont\bfseries Foundations of\\[0.3em]
Machine Learning\par}
\vspace{1.2cm}
\rule{0.55\textwidth}{0.8pt}\\[0.9em]
{\Large\itshape Lecture Notes by Dat Quoc Truong\par}
\vspace{0.9cm}
\rule{0.55\textwidth}{0.4pt}\\[1.2em]
{\large Chapters 1\,--\,10: From Linear Regression to
Reinforcement Learning\par}
\vfill
{\normalsize\itshape A Comprehensive Beginner's Textbook\par}
\vspace{1.5cm}
\end{titlepage}

\cleardoublepage

% ── Table of contents ─────────────────────────────────────
\tableofcontents
\clearpage

% ── Chapters ──────────────────────────────────────────────
% =========================================================
%  Chapter 1: Introduction to Machine Learning
% =========================================================
\chapter{Introduction to Machine Learning}

This chapter introduces the foundations of machine learning. We begin with a
high-level overview that motivates why machine learning matters and where it
appears in modern technology. We then distinguish supervised from unsupervised
learning, the two largest families of learning algorithms. After that, we focus
on the simplest yet most foundational supervised learning algorithm ---
\emph{linear regression with one variable} --- and study its model, its cost
function, and its visualisation. Finally, we examine \emph{gradient descent},
the optimisation algorithm used to train the model from data.

%----------------------------------------------------------
\section{Overview of Machine Learning}

\subsection{Why Machine Learning Matters}

Machine learning (ML) is a foundational technology behind countless applications
that shape modern life: web search engines, email spam filtering, speech
recognition, medical diagnosis, recommendation systems, autonomous driving, and
many more. Its strength lies in solving problems that are difficult or impossible
to express with explicit rules, by instead \emph{learning patterns directly from
data}.

Machine learning is widely considered a sub-field of artificial intelligence
(AI). It is sometimes confused with the more ambitious goal of
\emph{Artificial General Intelligence} (AGI), which refers to systems with
broad, human-level cognitive abilities. AGI remains largely hypothetical and is
often overhyped in popular media. Practical ML, in contrast, is a mature
engineering discipline with concrete tools and well-understood algorithms.

\subsection{What Is Machine Learning?}

The basic idea of machine learning is captured in a simple sentence:
\emph{learn to perform a task from data}. A more technical phrasing is that an
ML system \emph{recognises and extracts patterns} from data, and then uses those
patterns to make predictions, decisions, or descriptions about new, unseen
inputs.

A classical definition was given by Tom Mitchell (1997):

\begin{quote}
\itshape
A computer program is said to learn from experience $E$ with respect to some
class of tasks $T$ and performance measure $P$, if its performance on tasks in
$T$, as measured by $P$, improves with experience $E$.
\end{quote}

This definition highlights three key ingredients: \textbf{experience} (data the
system learns from), \textbf{tasks} (what the system must do), and
\textbf{performance measure} (how we evaluate success).

\subsection{Major Families of ML Algorithms}

Machine learning algorithms are commonly grouped into the following families:

\begin{itemize}
\item \textbf{Supervised learning.} The most widely used family in real-world
  applications. Each training example is a pair: an input $x$ and a
  corresponding correct output $y$. The algorithm learns a mapping
  $f\colon x \to y$.
\item \textbf{Unsupervised learning.} The data consists only of inputs $x$,
  with no labels. The goal is to discover hidden structure --- for example,
  clusters, low-dimensional patterns, or anomalies.
\item \textbf{Recommender systems.} A specialised family that predicts a
  user's preferences for items, used by streaming services, e-commerce
  platforms, and social media.
\item \textbf{Reinforcement learning.} An agent learns by interacting with an
  environment, receiving \emph{rewards} or \emph{penalties}, and gradually
  improving its behaviour to maximise cumulative reward.
\end{itemize}

\begin{figure}[h]
\centering
\begin{tikzpicture}[
  node distance=1.4cm and 2.6cm,
  box/.style={draw, rounded corners=4pt, minimum width=3.4cm,
              minimum height=0.9cm, font=\small, align=center,
              line width=0.7pt}]
\node[box, minimum width=4.0cm, line width=1pt, fill=gray!8]
  (ml) {\textbf{Machine Learning}};
\node[box, below left=1.5cm and 2.0cm of ml]  (sup)   {Supervised\\Learning};
\node[box, below=1.5cm of ml]                 (unsup) {Unsupervised\\Learning};
\node[box, below right=1.5cm and 2.0cm of ml] (rl)    {Reinforcement\\Learning};
\draw[-{Stealth[length=2.5mm]}] (ml.south west) -- (sup.north);
\draw[-{Stealth[length=2.5mm]}] (ml.south)      -- (unsup.north);
\draw[-{Stealth[length=2.5mm]}] (ml.south east) -- (rl.north);
\end{tikzpicture}
\caption{Three major families of machine learning algorithms.
Recommender systems are a specialised sub-area of supervised and unsupervised
learning.}
\label{fig:ml-families}
\end{figure}

%----------------------------------------------------------
\section{Supervised vs.\ Unsupervised Machine Learning}

\subsection{Supervised Learning}

In supervised learning, every training example has the form $(x,\,y)$, where
$x$ is the \emph{input} (also called a \emph{feature} or \emph{feature vector})
and $y$ is the correct \emph{output} (also called the \emph{label} or
\emph{target}). The objective is to learn a function $f$ such that
$f(x) \approx y$ on new, unseen inputs. The accuracy of the resulting model
depends on the difficulty of the task, the amount and quality of the dataset,
and the choice of algorithm.

Supervised learning problems split into two broad categories depending on the
type of the output variable.

\paragraph{Classification problems.}
The output $y$ is \emph{discrete}, i.e.\ a category label. Examples include:

\begin{itemize}
\item \textbf{Binary} (2 classes): yes/no, 0/1, true/false.
  \emph{Example: Is this email spam?}
\item \textbf{Multiclass} ($\ge3$ classes): classifying a fruit as
  \emph{apple}, \emph{mandarin}, or \emph{lemon} based on width and height.
\item \textbf{Multilabel}: an example may belong to several classes at once.
  \emph{Example: tagging a photo with both ``beach'' and ``sunset.''}
\end{itemize}

Common practical classification tasks include:
\begin{enumerate}
\item Detecting whether a new email is spam from its text content.
\item Predicting whether a user will buy a product, given features such as age,
  gender, and browsing history.
\item Forecasting whether tomorrow will be rainy, using humidity, wind speed,
  and cloudiness.
\item Detecting whether a credit-card transaction is fraudulent.
\end{enumerate}

\paragraph{Regression problems.}
The output $y$ is \emph{continuous} (a real number). The model predicts a real
value $\hat{y} = f(x)$, attempting to be as close to the true $y$ as possible.
Examples include predicting house price from size, temperature from time of day,
or a patient's blood pressure from lifestyle factors.

Figure~\ref{fig:cls-vs-reg} contrasts the two types visually.

\begin{figure}[h]
\centering
\begin{tikzpicture}
\begin{axis}[
  name=ax1,
  width=0.46\textwidth, height=5.5cm,
  title={\textbf{Classification}},
  xlabel={Feature $x_1$}, ylabel={Feature $x_2$},
  xtick=\empty, ytick=\empty,
  axis lines=left, enlargelimits=0.12,
  title style={font=\small\bfseries}]
\addplot[only marks, mark=*, mark size=2.2pt, color=pBlue]
  coordinates{(1,1)(1.5,2.1)(2,1.5)(2.2,2.5)(1.2,2.2)(1.8,1.1)};
\addplot[only marks, mark=square*, mark size=2pt, color=pRed]
  coordinates{(4,4)(4.5,3.5)(5,4.5)(3.8,4.2)(4.7,3.7)(4.2,4.8)};
\addplot[domain=0.5:6, samples=2, dashed, thick, pGray]{6-x};
\node[pBlue,  font=\small] at (axis cs:1.2,0.5) {Class 0};
\node[pRed,   font=\small] at (axis cs:4.8,5.0) {Class 1};
\end{axis}
\begin{axis}[
  name=ax2, at={(ax1.east)}, xshift=1.6cm,
  width=0.46\textwidth, height=5.5cm,
  title={\textbf{Regression}},
  xlabel={Input $x$}, ylabel={Target $y$},
  xtick=\empty, ytick=\empty,
  axis lines=left, enlargelimits=0.12,
  title style={font=\small\bfseries}]
\addplot[only marks, mark=*, mark size=2pt, color=pBlue]
  coordinates{(1,1.2)(2,1.8)(3,2.9)(4,3.2)(5,4.1)(6,4.7)(7,5.3)(8,6.1)};
\addplot[domain=0.5:8.5, samples=2, color=pRed, thick]{0.72*x+0.45};
\node[pRed, font=\small] at (axis cs:5.5,2.6) {$\hat{y}=f(x)$};
\end{axis}
\end{tikzpicture}
\caption{Left: classification separates discrete categories with a decision
boundary. Right: regression fits a continuous function to the data.}
\label{fig:cls-vs-reg}
\end{figure}

\subsection{Unsupervised Learning}

In unsupervised learning, the dataset consists only of inputs
$x^{(1)}, x^{(2)}, \dots, x^{(m)}$, with \emph{no} corresponding labels $y$.
The algorithm must discover useful structure entirely on its own, without any
guidance from human-provided answers.

Common types of unsupervised learning include:

\begin{itemize}
\item \textbf{Clustering.} Group similar data points together.
  \emph{Examples:} grouping news articles by topic; grouping customers by
  purchasing behaviour; analysing DNA microarrays to find gene patterns.
\item \textbf{Anomaly detection.} Find unusual data points that deviate from
  normal behaviour, e.g.\ unusual credit-card transactions for fraud monitoring.
\item \textbf{Dimensionality reduction.} Compress the data by representing it
  with fewer numbers while preserving important structure. Useful for
  visualisation and pre-processing.
\end{itemize}

\begin{figure}[h]
\centering
\begin{tikzpicture}
\begin{axis}[
  width=0.6\textwidth, height=5.5cm,
  title={\textbf{Clustering: discovering groups without labels}},
  title style={font=\small\bfseries},
  xlabel={Feature 1}, ylabel={Feature 2},
  xtick=\empty, ytick=\empty,
  axis lines=left, enlargelimits=0.1]
\addplot[only marks, mark=*, mark size=2.2pt, color=pBlue]
  coordinates{(1,1)(1.4,1.6)(1.8,1.2)(2.0,2.0)(1.3,2.2)(2.2,1.5)(0.9,1.8)};
\addplot[only marks, mark=*, mark size=2.2pt, color=pRed]
  coordinates{(5,5)(5.4,4.6)(5.8,5.2)(6.0,5.0)(5.3,5.5)(6.1,4.8)(5.7,5.8)};
\addplot[only marks, mark=*, mark size=2.2pt, color=pGreen]
  coordinates{(5,1)(5.4,1.6)(5.8,1.2)(6.0,2.0)(5.3,1.5)(6.1,1.8)(5.0,0.8)};
\addplot[only marks, mark=x, mark size=5pt, pBlue,  line width=1.2pt]
  coordinates{(1.51,1.7)};
\addplot[only marks, mark=x, mark size=5pt, pRed,   line width=1.2pt]
  coordinates{(5.61,5.13)};
\addplot[only marks, mark=x, mark size=5pt, pGreen, line width=1.2pt]
  coordinates{(5.51,1.56)};
\end{axis}
\end{tikzpicture}
\caption{A clustering algorithm assigns points to three groups (shown by colour)
without ever being told the correct group. Crosses mark the centroids.}
\label{fig:cluster}
\end{figure}

\subsection{Side-by-Side Comparison}

\begin{center}
\begin{tabular}{lp{5.2cm}p{5.2cm}}
\toprule
\textbf{Aspect} & \textbf{Supervised learning} & \textbf{Unsupervised learning}\\
\midrule
Training data  & Pairs $(x,\,y)$ with labels & Inputs $x$ only, no labels\\
Goal           & Learn $f$ s.t.\ $f(x)\approx y$ & Discover hidden structure\\
Typical tasks  & Classification, regression & Clustering, anomaly detection, dimensionality reduction\\
Classic example & Predicting house price from size & Grouping customers by behaviour\\
\bottomrule
\end{tabular}
\end{center}

%----------------------------------------------------------
\section{Linear Regression Model}

We now zoom in on the simplest and most important supervised learning model:
\textbf{linear regression with one variable}. Linear regression fits a
\emph{straight line} to a dataset and uses that line to predict numerical
outputs from a single numerical input.

\subsection{Setup and Notation}

Suppose we wish to predict house prices from house sizes. A typical training
set might look like:

\begin{center}
\begin{tabular}{cc}
\toprule
Size in ft$^2$ \quad ($x$, feature) & Price in \$1000s \quad ($y$, target)\\
\midrule
2104 & 460\\
1416 & 232\\
1534 & 315\\
\phantom{0}852 & 178\\
$\vdots$ & $\vdots$\\
\bottomrule
\end{tabular}
\end{center}

We adopt the following standard notation:

\begin{center}
\begin{tabular}{llp{7cm}}
\toprule
\textbf{Symbol} & \textbf{Name} & \textbf{Meaning}\\
\midrule
$m$         & Number of examples & Total rows in the dataset\\
$x$         & Input / feature    & E.g.\ house size in ft$^2$\\
$y$         & Output / target    & E.g.\ house price in \$1000s\\
$(x,\,y)$   & Training example   & A single (input, output) pair\\
$(x^{(i)},\,y^{(i)})$ & $i$-th example & Index $i$ is a row index, \emph{not} an exponent\\
$\hat{y}$   & Prediction         & Model output for a given $x$\\
\bottomrule
\end{tabular}
\end{center}

\subsection{The Linear Model}

Linear regression assumes the relationship between $x$ and $y$ is
approximately a straight line:
\begin{equation}
\hat{y} \;=\; f_{w,b}(x) \;=\; w\,x + b,
\label{eq:linmodel}
\end{equation}
where $w$ is the \emph{slope} (or \emph{weight}) and $b$ is the
\emph{intercept} (or \emph{bias}). Together, $w$ and $b$ are the
\textbf{parameters} of the model — the variables we adjust during training.
Because $y$ is a real number and the model's output is also a real number, this
is a \emph{regression} model (not classification).

\subsection{The Cost Function}

To choose good parameters $w$ and $b$ we need to measure how well the line
fits the data. This measure is called the \textbf{cost function} or
\emph{loss function}. For linear regression the standard choice is the
\textbf{mean squared error (MSE)} cost:
\begin{equation}
J(w,\,b) \;=\; \frac{1}{2m}\sum_{i=1}^{m}
\bigl(f_{w,b}(x^{(i)}) - y^{(i)}\bigr)^2
\;=\; \frac{1}{2m}\sum_{i=1}^{m}
\bigl(w\,x^{(i)} + b - y^{(i)}\bigr)^2.
\label{eq:mse-cost}
\end{equation}

The factor $\tfrac{1}{m}$ averages over all training examples, and the extra
$\tfrac{1}{2}$ is a convenience that makes differentiation cleaner (the $2$
from the power rule cancels). The \textbf{goal} is to choose $w,b$ to
\emph{minimise} $J(w,b)$: when $J$ is small, the model's predictions are
close to the true targets.

\subsection{Visualising the Cost Function}

To build intuition, temporarily fix $b=0$ so the model becomes
$f_w(x)=wx$. As we vary $w$, the line rotates around the origin, and
$J(w)$ traces a smooth, bowl-shaped (convex) curve.

\begin{figure}[h]
\centering
\begin{tikzpicture}
\begin{axis}[
  width=0.65\textwidth, height=6cm,
  xlabel={Parameter $w$}, ylabel={Cost $J(w)$},
  title={\textbf{Cost function is bowl-shaped (convex)}},
  title style={font=\small\bfseries},
  axis lines=left, samples=100, domain=-0.8:2.8,
  enlargelimits=0.06, xtick=\empty, ytick=\empty,
  every axis plot/.append style={line width=1.4pt}]
\addplot[pBlue]{(x-1)^2};
\addplot[only marks, mark=*, color=pRed, mark size=4pt]
  coordinates{(1,0)};
\node[above right, font=\small, pRed] at (axis cs:1.05,0) {global minimum};
\draw[dashed, pGray, thin] (axis cs:1,0) -- (axis cs:1,0.8);
\end{axis}
\end{tikzpicture}
\caption{For fixed $b$, the cost $J(w)$ is a convex (bowl-shaped) function with
a unique global minimum. This guarantees gradient descent will find the optimal
$w$.}
\label{fig:Jw}
\end{figure}

When both $w$ and $b$ are free, $J(w,b)$ is a bowl in three dimensions. We
can visualise it as a \emph{contour plot}: closed curves in the
$(w,b)$-plane where $J$ is constant.

\begin{figure}[h]
\centering
\begin{tikzpicture}
\begin{axis}[
  width=0.58\textwidth, height=5.6cm,
  xlabel={$w$}, ylabel={$b$},
  title={\textbf{Contour plot of $J(w,b)$}},
  title style={font=\small\bfseries},
  axis lines=left, enlargelimits=0.08,
  view={0}{90}, xtick=\empty, ytick=\empty,
  every axis plot/.append style={line width=0.9pt}]
\addplot[color=pBlue!90, domain=-pi:pi, samples=80]
  ({0.45*cos(deg(x))+1},{0.45*0.6*sin(deg(x))+1});
\addplot[color=pBlue!70, domain=-pi:pi, samples=80]
  ({0.9*cos(deg(x))+1},{0.9*0.6*sin(deg(x))+1});
\addplot[color=pBlue!50, domain=-pi:pi, samples=80]
  ({1.35*cos(deg(x))+1},{1.35*0.6*sin(deg(x))+1});
\addplot[color=pBlue!35, domain=-pi:pi, samples=80]
  ({1.8*cos(deg(x))+1},{1.8*0.6*sin(deg(x))+1});
\addplot[color=pBlue!22, domain=-pi:pi, samples=80]
  ({2.25*cos(deg(x))+1},{2.25*0.6*sin(deg(x))+1});
\addplot[only marks, mark=*, color=pRed, mark size=4pt] coordinates{(1,1)};
\node[above right, font=\small, pRed] at (axis cs:1.05,1) {minimum};
\end{axis}
\end{tikzpicture}
\caption{Contour plot of $J(w,b)$. Each elliptical curve is a level set
$J(w,b)=c$ for some constant $c$. The global minimum is at the centre.}
\label{fig:contour}
\end{figure}

\paragraph{Worked example.}
Consider the tiny training set $\{(1,1),(2,2),(3,3)\}$ with model
$f_w(x)=wx$ ($b=0$).

\begin{itemize}
\item $w=1$: predictions match targets exactly $\Rightarrow J(1)=0$.
\item $w=0.5$: predictions are $0.5,\,1.0,\,1.5$; squared errors are
  $0.25,\,1.0,\,2.25$; so $J(0.5)=\tfrac{3.5}{6}\approx0.583$.
\item $w=0$: predictions all zero; $J(0)=\tfrac{14}{6}\approx2.33$.
\end{itemize}

Among these three choices, $w=1$ is optimal — perfectly fitting the data.

%----------------------------------------------------------
\section{Training the Model with Gradient Descent}

We now have a clear objective: find $(w,b)$ that minimises $J(w,b)$. For linear
regression with one variable, a closed-form solution exists, but for almost all
other ML models we must rely on numerical optimisation. The most important such
algorithm is \textbf{gradient descent}.

\subsection{The Idea}

Imagine standing at some point on a hilly landscape that represents the cost
surface. Your goal is to reach the lowest valley. A sensible local strategy
is: \emph{look in every direction, find where the ground slopes most steeply
downward, and take a small step that way.} Repeat until you can no longer go
down. This is precisely what gradient descent does. At each iteration it
evaluates the \emph{gradient} (vector of partial derivatives) at the current
point and moves a small amount in the direction \emph{opposite} to the
gradient.

\paragraph{Algorithm outline.}
\begin{enumerate}
\item \textbf{Initialise} the parameters, typically $w=0$, $b=0$.
\item \textbf{Repeat until convergence:}
  update $w$ and $b$ in the direction that most reduces $J(w,b)$.
\item \textbf{Stop} when the parameters change very little between iterations
  (the algorithm has reached a local, or for convex costs, the global minimum).
\end{enumerate}

\subsection{The Update Rule}

The gradient descent update rule is applied \emph{simultaneously} to all
parameters at every iteration:
\begin{align}
w &\;\leftarrow\; w \;-\; \alpha\,\pd{J(w,b)}{w},\label{eq:gd-w}\\[4pt]
b &\;\leftarrow\; b \;-\; \alpha\,\pd{J(w,b)}{b}.\label{eq:gd-b}
\end{align}

Here $\alpha > 0$ is the \textbf{learning rate}, a small positive number that
controls the size of each step. The word \emph{simultaneous} is important: both
updates must use the \emph{old} values of $w$ and $b$. The correct implementation is:
\begin{align*}
\mathrm{tmp}_w &= w - \alpha\,\pd{J}{w}(w,b),\\
\mathrm{tmp}_b &= b - \alpha\,\pd{J}{b}(w,b),\\
w &\leftarrow \mathrm{tmp}_w,\quad b \leftarrow \mathrm{tmp}_b.
\end{align*}

\subsection{Derivative Intuition}

For a simplified one-variable setting ($b=0$), the update is
$w \leftarrow w - \alpha\,\tfrac{dJ}{dw}$.

\begin{itemize}
\item If the slope $\tfrac{dJ}{dw}>0$ (cost increases with $w$), the rule
  \emph{decreases} $w$ — moving toward the minimum.
\item If the slope $\tfrac{dJ}{dw}<0$ (cost decreases with $w$), the rule
  \emph{increases} $w$ — again moving toward the minimum.
\end{itemize}

\begin{figure}[h]
\centering
\begin{tikzpicture}
\begin{axis}[
  width=0.68\textwidth, height=5.8cm,
  xlabel={Parameter $w$}, ylabel={Cost $J(w)$},
  title={\textbf{Gradient descent descends toward the minimum}},
  title style={font=\small\bfseries},
  axis lines=left, samples=100, domain=-0.6:2.6,
  enlargelimits=0.06, xtick=\empty, ytick=\empty,
  every axis plot/.append style={line width=1.3pt}]
\addplot[pBlue]{(x-1)^2};
\addplot[only marks, mark=*, color=pRed, mark size=3.5pt]
  coordinates{(-0.4,2.0)(0.18,0.67)(0.57,0.19)(0.82,0.032)(1.0,0)};
\draw[-{Stealth[length=2.5mm]}, thick, pOrange]
  (axis cs:-0.4,2.0) -- (axis cs:0.18,0.67);
\draw[-{Stealth[length=2.5mm]}, thick, pOrange]
  (axis cs:0.18,0.67) -- (axis cs:0.57,0.19);
\draw[-{Stealth[length=2.5mm]}, thick, pOrange]
  (axis cs:0.57,0.19) -- (axis cs:0.82,0.032);
\draw[-{Stealth[length=2.5mm]}, thick, pOrange]
  (axis cs:0.82,0.032) -- (axis cs:1.0,0);
\node[above right, font=\small, pRed] at (axis cs:1.05,0) {minimum};
\end{axis}
\end{tikzpicture}
\caption{Successive gradient descent steps (orange arrows) descend toward the
minimum of the cost function.}
\label{fig:gd-steps}
\end{figure}

\subsection{Choosing the Learning Rate}

The learning rate $\alpha$ has a profound effect on training behaviour.

\textbf{Too small:} Each step is tiny, so convergence is very slow and may
require an impractical number of iterations to reach the minimum.

\textbf{Too large:} Steps may \emph{overshoot} the minimum, causing the cost
to oscillate or even grow --- the algorithm \emph{diverges}.

\begin{figure}[h]
\centering
\begin{tikzpicture}
\begin{axis}[
  name=axL,
  width=0.46\textwidth, height=5.2cm,
  title={$\alpha$ too small},
  title style={font=\small\bfseries},
  axis lines=left, samples=80, domain=-0.8:2.8,
  enlargelimits=0.08, xtick=\empty, ytick=\empty]
\addplot[pBlue, line width=1.2pt]{(x-1)^2};
\addplot[only marks, mark=*, color=pRed, mark size=2.5pt]
  coordinates{(-0.5,2.25)(-0.32,1.73)(-0.16,1.33)(0.0,1.0)(0.13,0.75)(0.24,0.57)};
\end{axis}
\begin{axis}[
  name=axR, at={(axL.east)}, xshift=1.6cm,
  width=0.46\textwidth, height=5.2cm,
  title={$\alpha$ too large (diverges)},
  title style={font=\small\bfseries},
  axis lines=left, samples=80, domain=-1.5:3.5,
  enlargelimits=0.08, xtick=\empty, ytick=\empty]
\addplot[pBlue, line width=1.2pt]{(x-1)^2};
\addplot[only marks, mark=*, color=pRed, mark size=2.5pt]
  coordinates{(-0.5,2.25)(2.6,2.56)(-1.0,4.0)(3.1,4.41)};
\draw[-{Stealth[length=2mm]}, pOrange, thick]
  (axis cs:-0.5,2.25)--(axis cs:2.6,2.56);
\draw[-{Stealth[length=2mm]}, pOrange, thick]
  (axis cs:2.6,2.56)--(axis cs:-1.0,4.0);
\draw[-{Stealth[length=2mm]}, pOrange, thick]
  (axis cs:-1.0,4.0)--(axis cs:3.1,4.41);
\end{axis}
\end{tikzpicture}
\caption{Left: a too-small $\alpha$ leads to painfully slow convergence.
Right: a too-large $\alpha$ causes overshooting and divergence.}
\label{fig:lr-too-small-large}
\end{figure}

\paragraph{Near a minimum.}
At a local minimum the gradient is zero, so the update term
$\alpha \cdot 0 = 0$ and parameters stop changing. Moreover, as the algorithm
approaches the minimum, the gradient magnitude \emph{decreases}, so steps
automatically become smaller even with a fixed $\alpha$. This means gradient
descent can converge with a fixed learning rate.

\subsection{Gradient Descent for Linear Regression}

Differentiating the MSE cost \eqref{eq:mse-cost} term by term:
\begin{align}
\pd{J}{w} &= \frac{1}{m}\sum_{i=1}^{m}
  \bigl(f_{w,b}(x^{(i)})-y^{(i)}\bigr)\,x^{(i)},\label{eq:dJdw}\\[6pt]
\pd{J}{b} &= \frac{1}{m}\sum_{i=1}^{m}
  \bigl(f_{w,b}(x^{(i)})-y^{(i)}\bigr).\label{eq:dJdb}
\end{align}

(The factor of 2 from differentiating the square cancels with the $\tfrac{1}{2}$
in front of the cost.) The full \textbf{batch gradient descent for linear
regression} algorithm is then:

\medskip
\textbf{Repeat until convergence:}
\begin{align*}
w &\;\leftarrow\; w - \frac{\alpha}{m}\sum_{i=1}^{m}
  \bigl(f_{w,b}(x^{(i)})-y^{(i)}\bigr)\,x^{(i)},\\[4pt]
b &\;\leftarrow\; b - \frac{\alpha}{m}\sum_{i=1}^{m}
  \bigl(f_{w,b}(x^{(i)})-y^{(i)}\bigr).
\end{align*}

\paragraph{Why ``batch''?}
This variant uses \emph{all} $m$ training examples to compute each gradient
update (all $m$ form a single ``batch''). Other variants --- stochastic gradient
descent, mini-batch gradient descent --- use one example or a small subset at
a time and are introduced in later chapters.

\paragraph{Convexity guarantee.}
For the squared-error linear regression cost, $J(w,b)$ is a \emph{convex}
function --- its 3D shape is a single bowl with no local minima other than the
global minimum. Consequently, gradient descent with a sufficiently small
$\alpha$ is \emph{guaranteed} to converge to the global minimum regardless of
where the parameters are initialised.

%----------------------------------------------------------
\section*{Chapter Summary}
\addcontentsline{toc}{section}{Chapter Summary}

\begin{itemize}
\item Machine learning is the discipline of automatically learning patterns from
  data, rather than hand-coding explicit rules.
\item The two largest families are \textbf{supervised learning} (labelled data,
  predict $y$ from $x$) and \textbf{unsupervised learning} (unlabelled data,
  discover hidden structure).
\item Supervised problems split into \textbf{classification} (discrete output)
  and \textbf{regression} (continuous output).
\item Linear regression models the relationship as a straight line
  $f_{w,b}(x)=wx+b$.
\item The \textbf{squared-error cost} $J(w,b)$ measures how well the line fits
  the data; the goal is to minimise $J$.
\item \textbf{Gradient descent} iteratively updates $w$ and $b$ by moving in
  the direction of steepest descent. The learning rate $\alpha$ controls the
  step size.
\item For squared-error linear regression, $J$ is convex, so gradient descent
  always reaches the global minimum.
\end{itemize}

%----------------------------------------------------------
\section*{Exercises}
\addcontentsline{toc}{section}{Exercises}

\begin{enumerate}
\item For each of the following tasks, state whether it is
  \emph{classification} or \emph{regression}, and justify your answer.
  \begin{enumerate}[label=(\alph*)]
  \item Predicting tomorrow's maximum temperature in Celsius.
  \item Recognising which handwritten digit (0--9) appears in an image.
  \item Estimating the monthly rental price of an apartment.
  \item Detecting whether a tumour is malignant or benign.
  \end{enumerate}
\item Given the dataset $\{(1,2),(2,4),(3,6)\}$ and the model $f_w(x)=wx$,
  compute $J(w)$ for $w\in\{1,\,1.5,\,2,\,2.5\}$.
  Which value of $w$ has the lowest cost? Does this make sense intuitively?
\item Suppose at iteration $t$ the gradient $\partial J/\partial w = -3$ and
  $\alpha=0.1$. By how much does $w$ change, and in which direction?
\item Explain in your own words why using a very large learning rate
  (e.g.\ $\alpha=100$) is problematic. What would the learning curve look like?
\item True or False: ``In supervised learning, the algorithm has to figure out
  which categories exist by itself.'' Justify your answer.
\item Sketch the cost function $J(w)$ for the dataset $\{(1,3),(2,5),(3,7)\}$
  and model $f_w(x)=wx$. Where is the minimum?
\end{enumerate}

% =========================================================
%  Chapter 2: Linear Regression with Multiple Variables
% =========================================================
\chapter{Linear Regression with Multiple Variables}

In Chapter~1 we studied simple linear regression with a \emph{single} input
feature $x$. Real-world prediction problems almost always involve
\emph{more than one} input variable. This chapter generalises the model so
that the target $y$ is expressed as a linear combination of $n$ features
$x_1, x_2, \ldots, x_n$. We call this \textbf{multiple linear regression}.

%----------------------------------------------------------
\section{Multiple Features}

\subsection{Motivation and Extended Notation}

\begin{example}[Housing price prediction with multiple features]
Suppose we want to predict the selling price of a house. Many attributes
influence the price simultaneously:

\begin{center}
\begin{tabular}{llcc}
\toprule
& \textbf{Feature} & \textbf{Description} & \textbf{Typical range}\\
\midrule
$x_1$ & Size (sq.\ ft.) & Living area & 300--2{,}000\\
$x_2$ & \# Bedrooms     & Count       & 1--5\\
$x_3$ & \# Floors       & Count       & 1--3\\
$x_4$ & Age (years)     & House age   & 0--80\\
\midrule
$y$   & Price (\$1{,}000s) & Target   & ---\\
\bottomrule
\end{tabular}
\end{center}
Each row of the training set is one house; each column (except $y$) is a
feature.
\end{example}

We extend our notation as follows.

\begin{definition}[Notation for multiple linear regression]
Let the training set contain $m$ examples, each described by $n$ features.
\begin{itemize}
\item $n$ --- total number of features.
\item $x_j$ --- the $j$-th feature $(j=1,\ldots,n)$.
\item $\vec{x}^{(i)}=\bigl[x_1^{(i)},x_2^{(i)},\ldots,x_n^{(i)}\bigr]^{\!\top}$
  --- the column vector of all feature values for example $i$.
\item $x_j^{(i)}$ --- the value of feature $j$ in the $i$-th example.
\item An overhead arrow ($\vec{\,\cdot\,}$) denotes a vector; bold font
  $\mathbf{M}$ denotes a matrix.
\end{itemize}
\end{definition}

\begin{remark}
In linear algebra, vector indexing conventionally starts at 1
(i.e.\ $w_1,w_2,\ldots$), whereas Python and NumPy use zero-based indexing
(\texttt{w[0]}, \texttt{w[1]}, \ldots). This distinction is important when
translating mathematical formulae into code. In particular,
\texttt{range(n)} iterates from $0$ to $n-1$, not from $1$ to $n$.
\end{remark}

\subsection{The Multiple Linear Regression Model}

\subsubsection{Scalar Form}

The model predicts $y$ as a weighted sum of all $n$ features plus a bias $b$:
\begin{equation}
f_{\vec{w},b}(\vec{x})
\;=\; w_1 x_1 + w_2 x_2 + \cdots + w_n x_n + b.
\label{eq:mlr-scalar}
\end{equation}
Each weight $w_j\in\R$ measures the \emph{marginal effect} of feature $x_j$:
a one-unit increase in $x_j$ changes the predicted output by $w_j$, holding
all other features fixed. The bias $b$ is the predicted output when all
features equal zero; it functions as a \emph{baseline}.

\subsubsection{Vector (Dot-Product) Form}

Collecting the weights into a row vector
$\vec{w}=[w_1,w_2,\ldots,w_n]$ and the features into a column vector
$\vec{x}=[x_1,x_2,\ldots,x_n]^{\top}$, Eq.~\eqref{eq:mlr-scalar} simplifies
to
\begin{equation}
\boxed{f_{\vec{w},b}(\vec{x})
\;=\; \vec{w}\cdot\vec{x}+b
\;=\; \sum_{j=1}^{n}w_j\,x_j+b.}
\label{eq:mlr-vec}
\end{equation}

\begin{example}[Interpreting learned parameters]
Suppose gradient descent yields
\[
f(\vec{x})=0.1\,x_1+4\,x_2-2\,x_3-0.5\,x_4+80,
\]
with features from Example~2.1 and prices in \$1{,}000s.
\begin{itemize}
\item $w_1=0.1$: each additional sq.\ ft.\ adds \$100 to the predicted price.
\item $w_2=4$: each extra bedroom adds \$4{,}000.
\item $w_3=-2$: each additional floor \emph{reduces} the price by \$2{,}000
  (e.g.\ reflecting higher maintenance).
\item $w_4=-0.5$: each additional year of age reduces the price by \$500.
\item $b=80$: baseline prediction of \$80{,}000.
\end{itemize}
\end{example}

\subsection{Vectorisation for Computational Efficiency}

\subsubsection{Why Vectorisation Matters}

Computing predictions via a \texttt{for} loop is \emph{sequential}: each
multiplication is executed one after another. Modern CPUs and GPUs contain
specialised hardware --- SIMD (Single Instruction, Multiple Data) units ---
that can perform many floating-point operations \emph{simultaneously}.
\textbf{Vectorisation} phrases computations as linear-algebra operations that
numerical libraries such as \textbf{NumPy} dispatch to that hardware
automatically. For models with thousands of features, the speedup can reduce
wall-clock time from hours to minutes.

\subsubsection{Three Coding Approaches Compared}

\begin{center}
\begin{tabular}{lll}
\toprule
\textbf{Method} & \textbf{Code snippet} & \textbf{Efficiency}\\
\midrule
Hardcoded      & \texttt{f = w[0]*x[0]+w[1]*x[1]+\ldots}  & Low — not scalable\\[2pt]
\texttt{for} loop & \texttt{for j in range(n): f += w[j]*x[j]} & Medium — readable but slow\\[2pt]
Vectorised     & \texttt{f = np.dot(w, x) + b}             & High — uses parallel hardware\\
\bottomrule
\end{tabular}
\end{center}

\subsubsection{NumPy Implementation}

\begin{lstlisting}[caption={Computing a prediction with NumPy}]
import numpy as np

# 4-feature housing example
w = np.array([0.1, 4.0, -2.0, -0.5])   # weight vector  (n,)
b = 80.0                                  # scalar bias
x = np.array([1500, 3, 2, 10])           # one training example (n,)

# Vectorised prediction (single BLAS call)
f = np.dot(w, x) + b
print(f"Predicted price: ${f:.1f}k")     # Predicted price: $228.0k
\end{lstlisting}

\subsubsection{Vectorised Gradient Descent Update}

Rather than looping over each weight, all updates are performed in a single
array operation:
\begin{lstlisting}[caption={Vectorised parameter update}]
# dw: NumPy array of partial derivatives, shape (n,)
# db: scalar partial derivative for b
alpha = 0.01

w = w - alpha * dw   # updates ALL weights simultaneously
b = b - alpha * db   # scalar bias update
\end{lstlisting}

\subsection{Feature Engineering}

Raw features are not always the most informative representation.
\textbf{Feature engineering} is the practice of constructing new features by
transforming or combining existing ones, guided by domain knowledge or
exploratory analysis.

\begin{example}[Land area as an engineered feature]
A housing dataset may contain $x_1$: lot frontage (width in metres) and
$x_2$: lot depth (metres). We engineer
\[
x_3 = x_1\times x_2 \quad\text{(total land area in m}^2\text{)}.
\]
If area is more predictive of price than either dimension individually, the
model will assign a large weight $w_3$. The expanded model is
$f(\vec{x})=w_1 x_1+w_2 x_2+w_3(x_1 x_2)+b$.
\end{example}

\subsection{Polynomial Regression}

Feature engineering is the bridge to \textbf{polynomial regression}: by
including powers or other non-linear transformations of the original features,
we fit \emph{curved} relationships while still using the linear regression
framework — because the model remains linear in the \emph{parameters}.

Let $x$ denote house size in sq.\ ft. Three natural candidates are:
\begin{align}
\text{Quadratic:}   &\quad f(x)=w_1 x+w_2 x^2+b,\label{eq:quad}\\
\text{Cubic:}       &\quad f(x)=w_1 x+w_2 x^2+w_3 x^3+b,\\
\text{Square root:} &\quad f(x)=w_1 x+w_2\sqrt{x}+b.
\end{align}
Each model treats $x,x^2,x^3,\sqrt{x}$ as \emph{separate features} fed into
the standard linear regression framework.

\begin{figure}[h]
\centering
\begin{tikzpicture}
\begin{axis}[
  width=11cm, height=6.5cm,
  xlabel={House size $x$ (sq.\ ft.)},
  ylabel={Predicted price (\$1{,}000s)},
  xmin=0, xmax=2100, ymin=0, ymax=540,
  grid=both, grid style={gray!15, thin},
  legend pos=north west, legend style={font=\small, fill=white},
  tick label style={font=\small}, label style={font=\small},
  every axis plot/.append style={line width=1.4pt}]
\addplot[black, dashed, domain=0:2000, samples=80]{0.15*x+50};
\addlegendentry{Linear: $0.15x+50$}
\addplot[pBlue, domain=100:2000, samples=120]{50+0.08*x+0.00004*(x*x)};
\addlegendentry{Quadratic}
\addplot[pRed, domain=100:2000, samples=120]
  {50+0.12*x-0.00003*(x*x)+0.000000015*(x*x*x)};
\addlegendentry{Cubic}
\addplot[pGreen, domain=1:2000, samples=150]{50+2.5*sqrt(x)};
\addlegendentry{Square root}
\end{axis}
\end{tikzpicture}
\caption{Four regression models on hypothetical housing data. The square-root
model (green) rises steeply then flattens, the cubic (red) gives a smooth
upward trend, and the quadratic (blue) eventually turns downward
(possibly unrealistic).}
\label{fig:poly}
\end{figure}

\textbf{Important:} Raising a feature to a high power dramatically expands its
numerical range:
\[
x\in[1,\;1{,}000],\quad
x^2\in[1,\;10^6],\quad
x^3\in[1,\;10^9].
\]
Without feature scaling, one polynomial term would dominate the gradient
update. The scaling techniques in Section~\ref{sec:feat_scale} must be applied
\emph{before} running gradient descent on polynomial features.

%----------------------------------------------------------
\section{Gradient Descent in Practice}

\subsection{Cost Function and Update Rules for Multiple Features}

With $m$ training examples and $n$ features, the cost function generalises
naturally:
\begin{equation}
J(\vec{w},b)=\frac{1}{2m}\sum_{i=1}^{m}
\Bigl(f_{\vec{w},b}(\vec{x}^{(i)})-y^{(i)}\Bigr)^2.
\label{eq:mlr-cost}
\end{equation}

Gradient descent updates all $n$ weights plus the bias \emph{simultaneously}:
\begin{align}
w_j &\leftarrow w_j-\alpha\,\pd{J}{w_j},\quad j=1,\ldots,n,\label{eq:mlr-gd-w}\\[4pt]
b   &\leftarrow b-\alpha\,\pd{J}{b},\label{eq:mlr-gd-b}
\end{align}
where the partial derivatives are:
\begin{align}
\pd{J}{w_j} &= \frac{1}{m}\sum_{i=1}^m
  \Bigl(f_{\vec{w},b}(\vec{x}^{(i)})-y^{(i)}\Bigr)x_j^{(i)},\label{eq:mlr-dJdw}\\[4pt]
\pd{J}{b}   &= \frac{1}{m}\sum_{i=1}^m
  \Bigl(f_{\vec{w},b}(\vec{x}^{(i)})-y^{(i)}\Bigr).\label{eq:mlr-dJdb}
\end{align}
Note that Eq.~\eqref{eq:mlr-dJdw} differs from its single-feature analogue
only by the extra factor $x_j^{(i)}$; the bias update \eqref{eq:mlr-dJdb}
is identical in both settings.

\subsubsection{Vectorised Implementation}

\begin{lstlisting}[caption={Vectorised gradient descent for multiple linear regression}]
import numpy as np

def compute_gradient(X, y, w, b):
    """
    X  : feature matrix (m, n)
    y  : target vector  (m,)
    w  : weight vector  (n,)
    b  : bias (scalar)
    Returns dw (n,), db (scalar)
    """
    m     = X.shape[0]
    y_hat = X @ w + b        # (m,)
    err   = y_hat - y        # (m,)
    dw    = (X.T @ err) / m  # (n,)
    db    = err.mean()       # scalar
    return dw, db

def gradient_descent(X, y, w, b, alpha, n_iter):
    for _ in range(n_iter):
        dw, db = compute_gradient(X, y, w, b)
        w = w - alpha * dw
        b = b - alpha * db
    return w, b
\end{lstlisting}

\subsubsection{Gradient Descent vs.\ the Normal Equation}

An alternative closed-form solution, the \textbf{Normal Equation}, finds the
optimal parameters in a single linear algebra step:
\[
\vec{w}^*=\bigl(X^{\top}X\bigr)^{-1}X^{\top}\mathbf{y}.
\]

\begin{center}
\begin{tabular}{lll}
\toprule
\textbf{Property} & \textbf{Gradient Descent} & \textbf{Normal Equation}\\
\midrule
Iterations        & Many (iterative) & None (one-shot)\\
Learning rate     & Must be tuned    & Not required\\
Complexity        & $O(k\cdot mn)$   & $O(n^3)$ for matrix inversion\\
Scalability       & Excellent ($n>10^6$ feasible) & Slow for $n>10{,}000$\\
Applicability     & Any ML algorithm & Linear regression only\\
\bottomrule
\end{tabular}
\end{center}

For large-scale problems, gradient descent (and its stochastic or mini-batch
variants) is the standard choice.

%----------------------------------------------------------
\subsection{Feature Scaling}
\label{sec:feat_scale}

\subsubsection{Motivation}

When features have very different numerical ranges, the cost function takes on
an elongated, elliptical shape in parameter space (Figure~\ref{fig:contours}).
Gradient descent then oscillates along the steep walls of these ellipses rather
than heading directly toward the minimum, leading to slow convergence. Scaling
features to comparable ranges makes the cost function nearly \emph{spherical}
and gradient descent converges in far fewer iterations.

\begin{figure}[h]
\centering
\begin{subfigure}[b]{0.46\linewidth}
\centering
\begin{tikzpicture}
\begin{axis}[
  width=6.4cm, height=5.5cm,
  xlabel={$w_1$}, ylabel={$w_2$},
  xmin=-2.2, xmax=2.2, ymin=-4.8, ymax=4.8,
  tick label style={font=\tiny}, label style={font=\small},
  title={\small\textbf{Unscaled features}},
  axis lines=left]
\foreach \rx/\ry in {0.3/2.4, 0.6/1.9, 1.0/1.4, 1.5/0.9}{
  \addplot[pBlue!50, domain=0:360, samples=80, no marks]
    ({\rx*cos(x)},{\ry*sin(x)});}
\addplot[black, very thick, -{Stealth[length=2.5mm]},
  mark=*, mark size=1.6pt, mark options={fill=black}]
  coordinates{(-1.4,3.5)(-0.7,-2.8)(-0.2,1.8)(0,-1.1)(0,0)};
\addplot[mark=*, mark size=3.5pt,
  mark options={fill=pRed, draw=pRed}] coordinates{(0,0)};
\end{axis}
\end{tikzpicture}
\caption{Oscillating, slow convergence}
\end{subfigure}
\hfill
\begin{subfigure}[b]{0.46\linewidth}
\centering
\begin{tikzpicture}
\begin{axis}[
  width=6.4cm, height=5.5cm,
  xlabel={$w_1$}, ylabel={$w_2$},
  xmin=-2.2, xmax=2.2, ymin=-2.2, ymax=2.2,
  tick label style={font=\tiny}, label style={font=\small},
  title={\small\textbf{Scaled features}},
  axis equal, axis lines=left]
\foreach \r in {0.4, 0.8, 1.2, 1.7}{
  \addplot[pBlue!50, domain=0:360, samples=80, no marks]
    ({\r*cos(x)},{\r*sin(x)});}
\addplot[black, very thick, -{Stealth[length=2.5mm]},
  mark=*, mark size=1.6pt, mark options={fill=black}]
  coordinates{(-1.5,1.4)(-0.8,0.7)(-0.2,0.2)(0,0)};
\addplot[mark=*, mark size=3.5pt,
  mark options={fill=pRed, draw=pRed}] coordinates{(0,0)};
\end{axis}
\end{tikzpicture}
\caption{Direct, rapid convergence}
\end{subfigure}
\caption{Contour plots of $J(w_1,w_2)$. Without scaling (left), elongated
ellipses cause a slow, oscillating gradient path. With scaling (right), the
contours are circular and gradient descent converges directly. Red dot marks
the global minimum.}
\label{fig:contours}
\end{figure}

\subsubsection{Feature Magnitude vs.\ Parameter Magnitude}

There is an \textbf{inverse relationship} between a feature's range and the
magnitude of its learned weight:
\begin{itemize}
\item A feature with a \emph{large} range (e.g.\ size $\approx 1{,}000$
  sq.\ ft.) needs only a \emph{small} weight ($w_1\approx 0.1$) to produce
  a reasonable contribution to the prediction.
\item A feature with a \emph{small} range (e.g.\ bedrooms $\approx 1$--$5$)
  needs a \emph{large} weight ($w_2\approx 50$) to have comparable influence.
\end{itemize}
Unequal weight magnitudes lead to cost-function contours with very different
curvature along each axis, causing the oscillatory behaviour seen above.

\subsubsection{Scaling Techniques}

\paragraph{A.\ Division by Maximum}
\begin{equation}
x_j^{\text{scaled}}=\frac{x_j}{\max(x_j)}.
\end{equation}
Maps all values to $[0,1]$ (assuming non-negative features). Simple but
sensitive to outliers.

\paragraph{B.\ Mean Normalisation}
\begin{equation}
x_j\;\leftarrow\;\frac{x_j-\mu_j}{\max(x_j)-\min(x_j)},
\end{equation}
where $\mu_j$ is the feature mean. Results fall approximately in $[-1,+1]$.

\paragraph{C.\ Z-Score Normalisation (Standardisation)}
\begin{equation}
x_j\;\leftarrow\;\frac{x_j-\mu_j}{\sigma_j},
\label{eq:zscore}
\end{equation}
where $\sigma_j$ is the standard deviation. After this transformation each
feature has \emph{zero mean} and \emph{unit variance}. This is the most widely
used method and is especially effective when features follow approximately
Gaussian distributions.

\begin{example}[Z-score normalisation in numbers]
Suppose $x_1$ (house size) has $\mu_1=1{,}000$ sq.\ ft.\ and $\sigma_1=500$.

\begin{center}
\begin{tabular}{lccc}
\toprule
House & Raw $x_1$ & Mean-normalised & Z-score\\
\midrule
A & $300$  & $\frac{300-1000}{1700}\approx-0.41$ & $\frac{300-1000}{500}=-1.40$\\[4pt]
B & $1000$ & $\frac{1000-1000}{1700}=0.00$        & $\frac{1000-1000}{500}=0.00$\\[4pt]
C & $2000$ & $\frac{2000-1000}{1700}\approx+0.59$ & $\frac{2000-1000}{500}=+2.00$\\
\bottomrule
\end{tabular}
\end{center}
\end{example}

\subsubsection{Practical Guidelines}

The target after scaling is approximately $-1\le x_j\le 1$ for each feature.

\begin{center}
\begin{tabular}{lll}
\toprule
\textbf{Feature range (before scaling)} & \textbf{Rescale?} & \textbf{Reason}\\
\midrule
$[-3,+3]$ or narrower           & Usually not & Acceptable range\\
$[-0.3,+0.3]$ or narrower       & Yes         & Too small; gradients negligible\\
$[-100,+100]$ or wider           & Yes         & Dominates gradient update\\
$[0,0.001]$ or very compressed  & Yes         & Creates near-zero steps\\
\bottomrule
\end{tabular}
\end{center}

\begin{remark}
Features need not be \emph{perfectly} comparable — merely \emph{roughly} in
the same ballpark. When in doubt, always scale: it virtually never hurts
performance and almost always accelerates convergence.
\end{remark}

\subsubsection{Implementation with NumPy and Scikit-learn}

\begin{lstlisting}[caption={Z-score normalisation in NumPy and Scikit-learn}]
import numpy as np
from sklearn.preprocessing import StandardScaler

# --- Manual (NumPy) ---
mu    = X_train.mean(axis=0)      # compute on TRAINING set only
sigma = X_train.std(axis=0)
X_train_sc = (X_train - mu) / sigma
X_test_sc  = (X_test  - mu) / sigma   # apply SAME stats to test set

# --- Scikit-learn ---
scaler = StandardScaler()
X_train_sc = scaler.fit_transform(X_train)   # fit + transform
X_test_sc  = scaler.transform(X_test)        # transform only (no re-fit!)
\end{lstlisting}

\begin{remark}
\textbf{Always} compute $\mu_j$ and $\sigma_j$ from the \emph{training set
only}, then apply those same constants to the validation and test sets.
Fitting the scaler on the test set constitutes \textbf{data leakage} and leads
to overly optimistic performance estimates.
\end{remark}

%----------------------------------------------------------
\subsection{Checking for Convergence}

\subsubsection{The Learning Curve}

A \textbf{learning curve} plots the training cost $J(\vec{w},b)$ against the
number of gradient descent iterations. It is the primary diagnostic tool for
verifying that training is proceeding correctly.

\begin{figure}[h]
\centering
\begin{tikzpicture}
\begin{axis}[
  width=10.5cm, height=6cm,
  xlabel={Iteration}, ylabel={Cost $J(\vec{w},b)$},
  xmin=0, xmax=500, ymin=0, ymax=1150,
  grid=both, grid style={gray!15, thin},
  tick label style={font=\small}, label style={font=\small},
  every axis plot/.append style={line width=1.4pt}]
\addplot[pBlue, domain=0:500, samples=200]{1050*exp(-0.016*x)+50};
\draw[pGray, dashed, thick] (axis cs:300,0)--(axis cs:300,1100);
\node[font=\small, fill=white, inner sep=2pt, align=center]
  at (axis cs:410,640) {Converged\\(curve flattens)};
\draw[-{Stealth[length=2mm]}, pGray, thick]
  (axis cs:388,580)--(axis cs:325,120);
\end{axis}
\end{tikzpicture}
\caption{A typical learning curve. The cost decreases rapidly at first, then
levels off as the algorithm converges. A curve that increases or oscillates
signals a learning rate that is too large or a bug in the code.}
\label{fig:learning-curve}
\end{figure}

Key observations:
\begin{itemize}
\item A correct implementation always produces a \emph{monotonically
  decreasing} curve.
\item The number of iterations needed varies enormously by problem — from
  as few as 30 to over 100{,}000.
\item When the curve flattens, gradient descent has approximately converged.
\item An \emph{increasing} or \emph{oscillating} curve is a red flag: either
  $\alpha$ is too large, or there is a bug (e.g.\ using $+\alpha$ instead of
  $-\alpha$).
\end{itemize}

\subsubsection{Automatic Convergence Test}

Declare convergence when the decrease in cost over one iteration is smaller
than a threshold $\varepsilon$ (e.g.\ $\varepsilon=10^{-3}$):
\[
\bigl|J^{(\text{iter})}-J^{(\text{iter}-1)}\bigr|\le\varepsilon.
\]
Choosing a reliable $\varepsilon$ is difficult in practice; visual inspection
of the learning curve is often more informative.

%----------------------------------------------------------
\subsection{Choosing the Learning Rate}

\subsubsection{Diagnosing a Poor Learning Rate}

\begin{itemize}
\item \textbf{Monotonically increasing cost}: $\alpha$ is almost certainly too
  large — the update step overshoots the minimum.
\item \textbf{Oscillating cost}: $\alpha$ is too large; parameters bounce from
  one side of the minimum to the other.
\item \textbf{Very slow decrease}: $\alpha$ is too small; convergence will take
  an impractical number of iterations.
\item \textbf{Cost increases even with tiny $\alpha$}: there is likely a
  \emph{code bug} — the most common error is using $+\alpha$ in the update rule.
\end{itemize}

\subsubsection{Debugging Strategy}

Set $\alpha$ to an extremely small value (e.g.\ $\alpha=10^{-10}$) and confirm
that $J$ decreases monotonically. If it does not, the bug lies in the gradient
computation, not the learning rate.

\subsubsection{Grid Search for $\alpha$}

Test a logarithmically spaced range, roughly tripling each candidate:
\[
0.001\;\to\;0.003\;\to\;0.01\;\to\;0.03\;\to\;0.1\;\to\;0.3\;\to\;1.
\]
For each candidate, run a small number of iterations and plot the learning
curve. Select the \textbf{largest} value of $\alpha$ that still produces a
consistently decreasing curve.

\begin{figure}[h]
\centering
\begin{tikzpicture}
\begin{axis}[
  width=10.5cm, height=6cm,
  xlabel={Iteration}, ylabel={Cost $J$},
  xmin=0, xmax=150, ymin=0, ymax=1150,
  grid=both, grid style={gray!15, thin},
  tick label style={font=\small}, label style={font=\small},
  legend pos=north east, legend style={font=\small, fill=white},
  every axis plot/.append style={line width=1.3pt}]
\addplot[pGray, domain=0:150, samples=100]{1050*exp(-0.002*x)+50};
\addlegendentry{$\alpha=0.001$}
\addplot[pBlue, domain=0:150, samples=100]{1050*exp(-0.020*x)+50};
\addlegendentry{$\alpha=0.01$}
\addplot[pRed, domain=0:150, samples=100]{1050*exp(-0.12*x)+50};
\addlegendentry{$\alpha=0.1$ (best)}
\end{axis}
\end{tikzpicture}
\caption{Learning curves for three candidate learning rates. The largest stable
value ($\alpha=0.1$) converges the fastest.}
\label{fig:lr-grid}
\end{figure}

%----------------------------------------------------------
\section*{Chapter Summary}
\addcontentsline{toc}{section}{Chapter Summary}

\begin{itemize}
\item \textbf{Multiple linear regression} models the target as
  $f_{\vec{w},b}(\vec{x})=\vec{w}\cdot\vec{x}+b$, a weighted sum of $n$
  features plus a bias.
\item Each weight $w_j$ captures the marginal effect of feature $x_j$,
  holding others fixed. Positive $w_j$ indicates a direct relationship;
  negative indicates an inverse relationship.
\item \textbf{Vectorisation} (via \texttt{np.dot}) replaces slow sequential
  loops with parallel hardware instructions, yielding dramatic speedups.
\item \textbf{Feature engineering} (e.g.\ products, ratios) and
  \textbf{polynomial regression} (powers of $x$) can reveal more predictive
  representations while preserving the linear regression framework.
\item \textbf{Feature scaling} (especially z-score normalisation) maps features
  to comparable ranges, converting elongated cost ellipses into nearly circular
  contours, greatly accelerating gradient descent.
\item The \textbf{learning curve} (cost vs.\ iterations) is the key diagnostic:
  a correct implementation yields a monotonically decreasing, eventually flat
  curve.
\item The \textbf{learning rate} $\alpha$ should be the largest value that
  produces stable, monotonic decrease; test candidates on a logarithmic grid.
\item The \textbf{Normal Equation} offers a closed-form alternative but is
  computationally impractical for $n\gtrsim 10{,}000$.
\end{itemize}

% =========================================================
%  Chapter 3: Classification
% =========================================================
\chapter{Classification}

In the previous two chapters we studied \emph{linear regression}, whose goal
was to predict a continuous numerical output such as a house price or a
temperature. We now make a fundamental shift: instead of predicting a number
along a continuous spectrum, we wish to assign each input to one of a
\emph{finite set of categories}. This task is called \textbf{classification}.

Classification problems arise everywhere in applied machine learning:
\begin{itemize}
\item \textbf{Spam detection}: is an incoming email spam or legitimate?
\item \textbf{Fraud detection}: is a credit-card transaction fraudulent?
\item \textbf{Medical diagnosis}: is a tumour malignant or benign?
\item \textbf{Image recognition}: which digit (0--9) appears in this image?
\item \textbf{Sentiment analysis}: is a product review positive or negative?
\end{itemize}

The simplest and most common form is \textbf{binary classification}, where the
output $y$ takes exactly two values. By convention we encode these as:
\[y\in\{0,\,1\},\]
where $y=1$ denotes the \emph{positive class} (presence of the condition) and
$y=0$ the \emph{negative class} (absence). This naming carries no moral
connotation.

%----------------------------------------------------------
\section{Classification with Logistic Regression}
\label{sec:logreg}

\subsection{Why Linear Regression Is Inadequate for Classification}

A natural first instinct is to reuse the linear regression model
$f_{\vec{w},b}(\vec{x})=\vec{w}\cdot\vec{x}+b$ for classification by
thresholding: predict $\hat{y}=1$ if $f>0.5$, else $\hat{y}=0$. While this
can work in simple situations, it suffers from two fundamental problems.

\paragraph{Problem 1: The outlier effect.}
Suppose we have a well-separated dataset of benign ($y=0$) and malignant
($y=1$) tumours plotted against tumour size. A linear model may yield a
reasonable boundary at some threshold $x^*$. Now add a single extreme outlier
— a very large tumour that is still malignant. The least-squares line is
pulled toward the outlier, shifting the boundary $x^*$ and misclassifying
previously correct points. Linear regression is globally sensitive to every
training point; one aberrant observation can corrupt the boundary for the rest.

\paragraph{Problem 2: Unbounded output range.}
Linear regression produces outputs in $(-\infty,+\infty)$. For a yes/no
problem, predicting values such as $1.8$ or $-0.4$ is awkward: they cannot be
interpreted as probabilities. We need a model whose output is
\emph{guaranteed} to lie in $[0,1]$.

\subsection{The Sigmoid (Logistic) Function}

The resolution is to \emph{squash} the linear output through a function that
maps $\R$ onto $(0,1)$. The standard choice is the \textbf{sigmoid function}:

\begin{definition}[Sigmoid Function]
\[
g(z)=\frac{1}{1+e^{-z}},\qquad z\in\R.
\]
\end{definition}

\paragraph{Key properties.}
\begin{itemize}
\item As $z\to+\infty$: $e^{-z}\to 0$, so $g(z)\to 1$.
\item As $z\to-\infty$: $e^{-z}\to\infty$, so $g(z)\to 0$.
\item At $z=0$: $g(0)=\tfrac{1}{2}$.
\item $g$ is smooth, strictly increasing, and S-shaped (sigmoidal).
\item Derivative: $g'(z)=g(z)\bigl(1-g(z)\bigr)$.
\end{itemize}

\begin{figure}[h]
\centering
\begin{tikzpicture}
\begin{axis}[
  width=10cm, height=6.5cm,
  xlabel={$z$}, ylabel={$g(z)$},
  xmin=-6.5, xmax=6.5, ymin=-0.05, ymax=1.1,
  xtick={-6,-4,-2,0,2,4,6},
  ytick={0,0.25,0.5,0.75,1},
  grid=both, grid style={gray!15, thin},
  tick label style={font=\small}, label style={font=\small},
  axis lines=left,
  every axis plot/.append style={line width=1.5pt}]
\addplot[pBlue, smooth, samples=200, domain=-6.5:6.5]
  {1/(1+exp(-x))};
\addplot[pGray, dashed, thin] coordinates{(-6.5,0.5)(6.5,0.5)};
\addplot[pGray, dashed, thin] coordinates{(0,-0.05)(0,0.5)};
\addplot[only marks, mark=*, color=pRed, mark size=3.5pt]
  coordinates{(0,0.5)};
\node[below right, font=\small] at (axis cs:0.1,0.49) {$(0,\;0.5)$};
\node[right, font=\small, pBlue] at (axis cs:3.5,0.92) {$g(z)\to 1$};
\node[right, font=\small, pBlue] at (axis cs:3.5,0.08) {$g(z)\to 0$};
\end{axis}
\end{tikzpicture}
\caption{The sigmoid (logistic) function. It maps every real number to $(0,1)$
and equals exactly $0.5$ at $z=0$.}
\label{fig:sigmoid}
\end{figure}

\subsection{The Logistic Regression Model}

Logistic regression is built in two stages.

\textbf{Stage 1 --- Linear combination.} Compute a weighted sum:
\[z=\vec{w}\cdot\vec{x}+b.\]

\textbf{Stage 2 --- Sigmoid activation.} Pass $z$ through the sigmoid:
\[
f_{\vec{w},b}(\vec{x})=g(z)=g(\vec{w}\cdot\vec{x}+b)
=\frac{1}{1+e^{-(\vec{w}\cdot\vec{x}+b)}}.
\]

The output $f_{\vec{w},b}(\vec{x})$ is interpreted as a \emph{probability}:
\[
f_{\vec{w},b}(\vec{x})=P(y=1\mid\vec{x};\;\vec{w},b).
\]
This is the estimated probability that the label equals 1 given the input
$\vec{x}$ and parameters $(\vec{w},b)$. The complementary probability is
$P(y=0\mid\vec{x})=1-f_{\vec{w},b}(\vec{x})$.

\begin{example}[Tumour classification]
Let $w=2$, $b=-5$, $x=3$ cm (tumour size). Then
\[z=2\cdot 3-5=1,\qquad g(1)=\frac{1}{1+e^{-1}}\approx 0.731.\]
The model reports a $73.1\%$ probability that the tumour is malignant.
\end{example}

\subsection{From Probability to Class Prediction}

To obtain a hard class label we introduce a \textbf{decision threshold}
$\tau$ (typically $\tau=0.5$):
\[
\hat{y}=\begin{cases}1&\text{if }f_{\vec{w},b}(\vec{x})\ge\tau,\\
0&\text{otherwise.}\end{cases}
\]
Because $g(z)\ge 0.5\iff z\ge 0$, the prediction rule is equivalent to:
\[\hat{y}=1\iff\vec{w}\cdot\vec{x}+b\ge 0.\]

\subsection{Decision Boundaries}

The \textbf{decision boundary} is the set of inputs $\vec{x}$ for which
$\vec{w}\cdot\vec{x}+b=0$ (i.e.\ where $f=0.5$, the model is exactly neutral).

\paragraph{Linear boundaries.}
With two features $x_1,x_2$ and parameters $w_1,w_2,b$, the boundary is the
straight line $w_1 x_1+w_2 x_2+b=0$.
\textit{Example:} $w_1=1$, $w_2=1$, $b=-3$ gives boundary $x_1+x_2=3$.
Points with $x_1+x_2>3$ are classified as $y=1$.

\paragraph{Non-linear boundaries.}
By augmenting the feature vector with polynomial terms, the boundary in the
\emph{original} feature space can be non-linear. With features
$(x_1^2,x_2^2)$ and $w_1=w_2=1$, $b=-1$, the boundary becomes
$x_1^2+x_2^2=1$, a unit circle.

\begin{figure}[h]
\centering
\begin{tikzpicture}
\begin{axis}[
  width=7.5cm, height=7.5cm,
  xlabel={$x_1$}, ylabel={$x_2$},
  xmin=-2.5, xmax=2.5, ymin=-2.5, ymax=2.5,
  xtick={-2,-1,0,1,2}, ytick={-2,-1,0,1,2},
  grid=both, grid style={gray!15, thin},
  tick label style={font=\small}, label style={font=\small},
  axis lines=center]
\addplot[pBlue, thick, smooth, domain=0:360, samples=200]
  ({cos(x)},{sin(x)});
\addplot[only marks, mark=*, color=pBlue, mark size=2pt]
  coordinates{(0.3,0.3)(0.4,-0.3)(-0.3,0.4)(-0.2,-0.2)(0.0,0.5)};
\addplot[only marks, mark=square*, color=pRed, mark size=2pt]
  coordinates{(1.5,1.2)(1.3,-1.4)(-1.6,0.8)(-1.2,-1.5)(1.8,-0.5)
              (0.5,1.8)(-0.4,2.0)(2.0,0.4)};
\node[pBlue,  font=\small] at (axis cs:0,0)    {$\hat{y}=0$};
\node[pRed,   font=\small] at (axis cs:1.9,1.9) {$\hat{y}=1$};
\end{axis}
\end{tikzpicture}
\caption{Circular decision boundary arising from polynomial features
$(x_1^2,x_2^2)$ with $w_1=w_2=1$, $b=-1$.}
\label{fig:circle-boundary}
\end{figure}

%----------------------------------------------------------
\section{Cost Function for Logistic Regression}
\label{sec:logcost}

\subsection{Why Squared Error Fails}

Reusing the MSE cost from linear regression with the non-linear sigmoid output
produces a \emph{non-convex} cost surface with multiple local minima and flat
plateaus. Gradient descent may converge to a sub-optimal solution. We need a
cost function that (a) preserves convexity and (b) strongly penalises confident
wrong predictions.

\subsection{Loss vs.\ Cost: A Key Distinction}

\begin{definition}[Loss and Cost]
The \textbf{loss} $L\bigl(f,y\bigr)$ measures the penalty on a \emph{single}
training example: the cost of predicting $f$ when the true label is $y$.
The \textbf{cost} $J$ is the \emph{average loss} over all $m$ training examples:
\[J(\vec{w},b)=\frac{1}{m}\sum_{i=1}^m L\!\bigl(f_{\vec{w},b}(\vec{x}^{(i)}),y^{(i)}\bigr).\]
\end{definition}

\subsection{The Logistic (Cross-Entropy) Loss}

We define the loss piecewise:
\begin{align}
\text{If }y=1:&\quad L=-\log\!\bigl(f\bigr),\label{eq:loss-y1}\\
\text{If }y=0:&\quad L=-\log\!\bigl(1-f\bigr).\label{eq:loss-y0}
\end{align}

\paragraph{Interpretation for $y=1$ (Eq.~\ref{eq:loss-y1}).}
\begin{itemize}
\item When $f\approx 1$ (correctly confident): $-\log(1)=0$ — no penalty.
\item When $f\to 0$ (confidently wrong): $-\log(f)\to+\infty$ — catastrophic penalty.
\end{itemize}

\paragraph{Interpretation for $y=0$ (Eq.~\ref{eq:loss-y0}).}
\begin{itemize}
\item When $f\approx 0$ (correct): $-\log(1)=0$ — no penalty.
\item When $f\to 1$ (confidently wrong): $-\log(0)\to+\infty$ — catastrophic penalty.
\end{itemize}

\begin{figure}[h]
\centering
\begin{tikzpicture}
\begin{axis}[
  width=11cm, height=6.5cm,
  xlabel={Prediction $f$}, ylabel={Loss $L$},
  xmin=0.01, xmax=0.99, ymin=0, ymax=5,
  xtick={0,0.2,0.4,0.6,0.8,1},
  ytick={0,1,2,3,4,5},
  grid=both, grid style={gray!15, thin},
  tick label style={font=\small}, label style={font=\small},
  axis lines=left,
  legend style={at={(0.5,0.97)}, anchor=north, font=\small, fill=white},
  every axis plot/.append style={line width=1.4pt}]
\addplot[pBlue, smooth, samples=200, domain=0.005:0.995]{-ln(x)};
\addlegendentry{$L=-\log(f)$\quad ($y=1$)}
\addplot[pRed, dashed, smooth, samples=200, domain=0.005:0.995]{-ln(1-x)};
\addlegendentry{$L=-\log(1-f)$\quad ($y=0$)}
\end{axis}
\end{tikzpicture}
\caption{Logistic loss functions for $y=1$ (solid, blue) and $y=0$ (dashed,
red). Both approach $+\infty$ when the prediction is maximally wrong.}
\label{fig:loss-curves}
\end{figure}

\subsection{The Unified Loss Formula}

The piecewise definition collapses elegantly into a single expression:
\begin{equation}
\boxed{L\bigl(f,y\bigr)
=-y\,\log(f)-(1-y)\,\log(1-f).}
\label{eq:unified-loss}
\end{equation}

\paragraph{Verification.}
\begin{itemize}
\item If $y=1$: $(1-y)=0$ annihilates the second term $\Rightarrow L=-\log(f)$.\checkmark
\item If $y=0$: $y=0$ annihilates the first term $\Rightarrow L=-\log(1-f)$.\checkmark
\end{itemize}

\subsection{The Full Logistic Regression Cost Function}

Averaging the unified loss over all $m$ training examples:
\begin{equation}
\boxed{J(\vec{w},b)=-\frac{1}{m}\sum_{i=1}^m\Bigl[
y^{(i)}\log\!\bigl(f_{\vec{w},b}(\vec{x}^{(i)})\bigr)
+(1-y^{(i)})\log\!\bigl(1-f_{\vec{w},b}(\vec{x}^{(i)})\bigr)\Bigr].}
\label{eq:logistic-cost}
\end{equation}

This is the \textbf{cross-entropy loss} or \textbf{log-loss}. It has two
complementary justifications:
\begin{enumerate}
\item \textbf{Statistical (MLE).} Eq.~\eqref{eq:logistic-cost} is exactly the
  negative log-likelihood of the training data under a Bernoulli model with
  success probability $f^{(i)}$. Minimising cross-entropy is equivalent to
  Maximum Likelihood Estimation — a principled statistical objective.
\item \textbf{Convexity.} Because the sigmoid is log-concave, this cost
  function is \emph{convex} in $(\vec{w},b)$. The cost surface has a
  \emph{single global minimum} and no local minima, so gradient descent is
  guaranteed to converge to the optimal solution regardless of initialisation.
\end{enumerate}

%----------------------------------------------------------
\section{Gradient Descent for Logistic Regression}
\label{sec:loggd}

\subsection{The Training Objective}

We seek $(\vec{w}^*,b^*)=\argmin_{\vec{w},b}\;J(\vec{w},b)$.
Once found, the trained model estimates
$P(y=1\mid\vec{x};\;\vec{w}^*,b^*)$ for any new input $\vec{x}$.

\subsection{The Update Rules}

Gradient descent applies simultaneously:
\begin{align}
w_j&\leftarrow w_j-\alpha\,\pd{J}{w_j},\quad j=1,\ldots,n,\\[4pt]
b  &\leftarrow b-\alpha\,\pd{J}{b}.
\end{align}

Differentiating Eq.~\eqref{eq:logistic-cost} (using the chain rule and the
identity $g'(z)=g(z)(1-g(z))$):
\begin{align}
\pd{J}{w_j}&=\frac{1}{m}\sum_{i=1}^m
\bigl(f_{\vec{w},b}(\vec{x}^{(i)})-y^{(i)}\bigr)\,x_j^{(i)},\label{eq:log-dJdw}\\[4pt]
\pd{J}{b}  &=\frac{1}{m}\sum_{i=1}^m
\bigl(f_{\vec{w},b}(\vec{x}^{(i)})-y^{(i)}\bigr).\label{eq:log-dJdb}
\end{align}

\begin{remark}
The gradient expressions~\eqref{eq:log-dJdw}--\eqref{eq:log-dJdb} are
\emph{structurally identical} to those of linear regression. The critical
difference lies in the definition of $f$: in linear regression
$f=\vec{w}\cdot\vec{x}+b$; in logistic regression
$f=g(\vec{w}\cdot\vec{x}+b)$. Same form, different meaning.
\end{remark}

\paragraph{Derivation sketch for $\partial J/\partial w_j$.}
Let $f^{(i)}=g(z^{(i)})$, $z^{(i)}=\vec{w}\cdot\vec{x}^{(i)}+b$.
The single-example loss is
$\ell^{(i)}=-y^{(i)}\log f^{(i)}-(1-y^{(i)})\log(1-f^{(i)})$.
Using $\partial f^{(i)}/\partial w_j=f^{(i)}(1-f^{(i)})x_j^{(i)}$:
\[
\frac{\partial\ell^{(i)}}{\partial w_j}
=\Bigl(-\frac{y^{(i)}}{f^{(i)}}+\frac{1-y^{(i)}}{1-f^{(i)}}\Bigr)
f^{(i)}(1-f^{(i)})x_j^{(i)}
=\bigl(f^{(i)}-y^{(i)}\bigr)x_j^{(i)}.
\]
Averaging over all $m$ examples gives Eq.~\eqref{eq:log-dJdw}.\quad$\square$

\subsection{Practical Considerations}

\paragraph{Choosing the learning rate $\alpha$.}
Monitor the learning curve (cost vs.\ iteration). A well-chosen $\alpha$
produces a monotonically decreasing, eventually flat curve.

\paragraph{Feature scaling.}
When features have vastly different magnitudes, the cost surface becomes an
elongated ellipse. Standardising features (zero mean, unit variance) transforms
the contours to circles, allowing much faster convergence.

\paragraph{Vectorisation.}
For large datasets, express the gradient update as matrix--vector operations
(NumPy) to exploit parallelism.

\paragraph{Convergence criterion.}
Halt when the relative change in $J$ between successive iterations falls below
a tolerance $\varepsilon$:
\[\frac{|J^{(t)}-J^{(t-1)}|}{|J^{(t-1)}|+\varepsilon_0}<\varepsilon.\]

%----------------------------------------------------------
\section{The Problem of Overfitting}
\label{sec:overfit}

\subsection{The Bias--Variance Trade-off}

Every supervised learning algorithm navigates a fundamental tension: a model
that is too simple cannot capture the true structure (\textbf{underfitting},
or \emph{high bias}), while a model that is too complex memorises the training
data and fails to generalise (\textbf{overfitting}, or \emph{high variance}).

\paragraph{Underfitting (High Bias).}
An underfitting model has a strong preconception about the form of the
relationship. Fitting a straight line to clearly curved data forces the model
to ignore obvious patterns. Such a model performs poorly even on the training
set.

\paragraph{Overfitting (High Variance).}
An overfitting model is excessively complex — perhaps a degree-15 polynomial
fitted to 20 data points. It interpolates the training data perfectly (near-zero
training error) but its oscillatory curve is driven by noise rather than the
true signal. Small changes in the training set produce drastically different
models.

\paragraph{Good generalisation.}
The ideal model captures the true underlying pattern while ignoring random noise.
It achieves low error on both training and unseen test data.

\begin{figure}[h]
\centering
\begin{tikzpicture}
\begin{axis}[
  name=ax1,
  width=4.8cm, height=4.6cm, title={\small\textbf{Underfitting}},
  title style={font=\small\bfseries},
  xmin=0, xmax=4, ymin=0, ymax=4,
  xtick=\empty, ytick=\empty, axis lines=left]
\addplot[only marks, mark=*, mark size=1.8pt, pBlue]
  coordinates{(0.3,0.4)(0.7,0.9)(1.1,1.5)(1.5,2.0)(2.0,2.3)
              (2.5,2.0)(3.0,1.5)(3.5,0.8)(3.9,0.3)};
\addplot[pRed, thick, domain=0:4]{1.2};
\end{axis}
\begin{axis}[
  name=ax2, at={(ax1.east)}, xshift=1.4cm,
  width=4.8cm, height=4.6cm, title={\small\textbf{Good generalisation}},
  title style={font=\small\bfseries},
  xmin=0, xmax=4, ymin=0, ymax=4,
  xtick=\empty, ytick=\empty, axis lines=left]
\addplot[only marks, mark=*, mark size=1.8pt, pBlue]
  coordinates{(0.3,0.4)(0.7,0.9)(1.1,1.5)(1.5,2.0)(2.0,2.3)
              (2.5,2.0)(3.0,1.5)(3.5,0.8)(3.9,0.3)};
\addplot[pRed, thick, smooth, domain=0.1:3.9, samples=80]
  {-0.57*(x-2)^2+2.3};
\end{axis}
\begin{axis}[
  name=ax3, at={(ax2.east)}, xshift=1.4cm,
  width=4.8cm, height=4.6cm, title={\small\textbf{Overfitting}},
  title style={font=\small\bfseries},
  xmin=0, xmax=4, ymin=-0.5, ymax=4.5,
  xtick=\empty, ytick=\empty, axis lines=left]
\addplot[only marks, mark=*, mark size=1.8pt, pBlue]
  coordinates{(0.3,0.4)(0.7,0.9)(1.1,1.5)(1.5,2.0)(2.0,2.3)
              (2.5,2.0)(3.0,1.5)(3.5,0.8)(3.9,0.3)};
\addplot[pRed, thick, smooth, domain=0.2:3.95, samples=200]
  {-0.57*(x-2)^2+2.3+0.62*sin(deg(5.3*x))};
\end{axis}
\end{tikzpicture}
\caption{Three models fitted to the same data. \textit{Left}: underfitting (too
simple). \textit{Centre}: good generalisation. \textit{Right}: overfitting (too
complex, fitting the noise).}
\label{fig:fitting}
\end{figure}

\subsection{Remedies for Overfitting}

\subsubsection{Strategy 1: Collect More Training Data}

The most direct remedy. With more data, any fixed-complexity model is less able
to memorise individual points and is forced to learn the true underlying pattern.
Unfortunately, collecting more data is not always feasible — it may be
expensive, time-consuming, or impossible for rare events.

\subsubsection{Strategy 2: Feature Selection}

If the number of features $n$ is large relative to $m$, overfitting is likely.
Reducing the feature set — keeping only the most informative features — reduces
the model's capacity to overfit. Feature selection can be done manually
(guided by domain knowledge) or automatically (using statistical tests, embedded
methods, or regularisation paths).

\subsubsection{Strategy 3: Regularisation}
\label{sec:regularisation}

Regularisation is generally the preferred remedy because it allows us to retain
\emph{all} features while preventing any single parameter from dominating the
model.

\paragraph{Core idea.}
Large weight parameters are what allow a model to oscillate wildly. By
\emph{penalising} large weights in the cost function, the optimiser is forced
to keep them small, producing a smoother, more parsimonious model.

\paragraph{Regularised cost function ($L_2$ / Ridge).}
\begin{equation}
J_{\text{reg}}(\vec{w},b)
=J(\vec{w},b)+\frac{\lambda}{2m}\sum_{j=1}^n w_j^2,
\label{eq:reg-cost}
\end{equation}
where $\lambda\ge 0$ is the \textbf{regularisation parameter}. Note that $b$
is conventionally \emph{not} regularised.

\paragraph{The role of $\lambda$.}
\begin{itemize}
\item $\lambda=0$: no regularisation; the model may overfit.
\item $\lambda\to\infty$: all weights forced to zero; model underfits
  (predicts constant $f=\sigma(b)$).
\item $\lambda$ ``just right'': balance between fitting data and model
  simplicity, achieving good generalisation.
\end{itemize}

\subsection{Regularised Linear Regression}

\paragraph{Gradient descent updates.}
\begin{align}
w_j&\leftarrow w_j-\alpha\Bigl[\frac{1}{m}\sum_{i=1}^m
  \bigl(f_{\vec{w},b}(\vec{x}^{(i)})-y^{(i)}\bigr)x_j^{(i)}
  +\frac{\lambda}{m}w_j\Bigr],\label{eq:reg-lin-w}\\
b  &\leftarrow b-\frac{\alpha}{m}\sum_{i=1}^m
  \bigl(f_{\vec{w},b}(\vec{x}^{(i)})-y^{(i)}\bigr).\label{eq:reg-lin-b}
\end{align}

\paragraph{Weight decay interpretation.}
Rearranging \eqref{eq:reg-lin-w}:
\[
w_j\leftarrow\underbrace{\Bigl(1-\frac{\alpha\lambda}{m}\Bigr)}_{\text{weight decay factor}}w_j
-\frac{\alpha}{m}\sum_{i=1}^m\bigl(f^{(i)}-y^{(i)}\bigr)x_j^{(i)}.
\]
Since $\alpha\lambda/m$ is a small positive number (e.g.\ $0.0002$), the weight
decay factor is slightly less than 1. At every step, each weight is
\emph{gently shrunk} before the gradient update, creating a constant restoring
force toward zero.

\subsection{Regularised Logistic Regression}

The same approach applies directly:
\begin{equation}
J(\vec{w},b)=-\frac{1}{m}\sum_{i=1}^m\!\Bigl[
y^{(i)}\log f^{(i)}+(1-y^{(i)})\log(1-f^{(i)})\Bigr]
+\frac{\lambda}{2m}\sum_{j=1}^n w_j^2,
\end{equation}
with gradient descent updates identical in form to
\eqref{eq:reg-lin-w}--\eqref{eq:reg-lin-b} but with the sigmoid $f$.

\paragraph{Summary table.}
\begin{center}
\begin{tabular}{llll}
\toprule
\textbf{Situation} & \textbf{$J_{\text{train}}$} &
\textbf{$J_{\text{test}}$} & \textbf{Remedy}\\
\midrule
Underfitting (High Bias)     & High  & High  & More features; reduce $\lambda$\\
Good generalisation          & Low   & Low   & ---\\
Overfitting (High Variance)  & Low   & High  & More data; reduce features; increase $\lambda$\\
\bottomrule
\end{tabular}
\end{center}

%----------------------------------------------------------
\section*{Chapter Summary}
\addcontentsline{toc}{section}{Chapter Summary}

\begin{itemize}
\item \textbf{Logistic regression} uses the sigmoid function to map a linear
  combination of features onto $(0,1)$, yielding a probability estimate
  $P(y=1\mid\vec{x})$.
\item The decision boundary $\vec{w}\cdot\vec{x}+b=0$ is linear; polynomial
  features extend it to non-linear shapes (circles, ellipses, etc.).
\item The \textbf{cross-entropy cost} is derived from MLE and is convex,
  guaranteeing a unique global minimum. Squared error would produce a
  non-convex cost with local minima.
\item The gradient formulas for logistic regression look identical to those of
  linear regression, but $f$ is the sigmoid in the logistic case.
\item \textbf{Overfitting} occurs when the model is too complex relative to the
  training data. The three main remedies are more data, feature selection, and
  $L_2$ \textbf{regularisation} (controlled by $\lambda$).
\item Regularisation adds a weight-penalty term to the cost, shrinking weights
  at every gradient step (``weight decay''), smoothing the decision boundary.
\end{itemize}

% =========================================================
%  Chapter 4: Neural Networks
% =========================================================
\chapter{Neural Networks}

%----------------------------------------------------------
\section{Neural Networks Intuition}

\subsection{Biological Inspiration and Engineering Reality}

The original motivation behind \textbf{Artificial Neural Networks (ANNs)} ---
often called simply \emph{neural networks} or \emph{deep learning} models ---
was to create software that mimics the learning processes of the human brain.
While modern engineering has largely moved beyond strict biological analogies,
the foundational vocabulary still draws on neuroscience.

A \textbf{biological neuron} processes information in three stages:
\begin{enumerate}
\item \textbf{Dendrites} receive incoming electrical impulses from other neurons.
\item The \textbf{cell body (soma)} aggregates and processes these signals.
\item The \textbf{axon} transmits the resulting output to downstream neurons.
\end{enumerate}

An \textbf{artificial neuron} is a mathematical abstraction of this process.
It takes numerical inputs, computes a weighted sum, applies a non-linear
function, and emits an \textit{activation value}.

\begin{figure}[h]
\centering
\begin{tikzpicture}[>=Stealth, thick, font=\small]
% Biological neuron
\node[draw, ellipse, fill=nodeGray!50, minimum width=2.4cm,
      minimum height=1.3cm] (soma) at (0,0) {Cell Body};
\foreach \y/\lbl in {0.7/Dendrite 1, 0/{Dendrite 2}, -0.7/{Dendrite 3}}{
  \draw[->] (-3.0,\y) node[left]{\lbl} -- (soma);}
\draw[->] (soma) -- (3.0,0) node[right]{Axon};
% separator
\draw[dashed, pGray] (4.2,1.3) -- (4.2,-1.3);
\node at (4.2,1.55) {\small\textit{vs.}};
% Artificial neuron
\node[draw, circle, fill=nodeBlue!60, minimum size=1.6cm]
     (anode) at (7.5,0) {$\Sigma\;g$};
\foreach \y/\lbl in {0.7/{$x_1$}, 0/{$x_2$}, -0.7/{$x_3$}}{
  \draw[->] (5.4,\y) node[left]{\lbl} -- (anode);}
\draw[->] (anode) -- (9.6,0) node[right]{$a$};
% labels
\node[above=0.35cm of soma,  font=\footnotesize\itshape]
     {Biological neuron};
\node[above=0.35cm of anode, font=\footnotesize\itshape]
     {Artificial neuron};
\end{tikzpicture}
\caption{Structural analogy between a biological neuron and an artificial
neuron. The artificial neuron computes a weighted sum ($\Sigma$) and applies a
non-linear function ($g$) to produce the activation $a$.}
\label{fig:bio-neuron}
\end{figure}

\begin{remark}
Despite the name, modern neural networks have very little to do with how the
brain actually works. We still understand remarkably little about neuroscience.
Today's AI systems are driven primarily by \textit{engineering principles} and
large-scale empirical results rather than by neurological simulation.
\end{remark}

\subsection{A Brief History and the Data-Compute Story}

The development of neural networks has passed through alternating periods of
enthusiasm and neglect:

\begin{center}
\begin{tabular}{lll}
\toprule
\textbf{Era} & \textbf{Status} & \textbf{Key developments}\\
\midrule
1950s & Inception & Brain-inspired computing first proposed\\
1980s--early 1990s & First wave & Handwritten digit recognition (zip codes)\\
Late 1990s & Decline & Outperformed by SVMs and other ML methods\\
2005--present & Resurgence & Rebranded as ``Deep Learning''; breakthroughs\\
\bottomrule
\end{tabular}
\end{center}

The key milestones in the modern era include: speech recognition (early 2010s),
ImageNet 2012 (landmark computer vision result), NLP with Transformers, and
today's large language models.

\paragraph{Why now?}
Traditional algorithms exhibit a characteristic performance ceiling; beyond a
certain amount of training data, adding more examples yields diminishing
returns. Large neural networks, by contrast, continue to improve as datasets
grow. Two enabling factors: the explosion of digital data from the internet,
and GPU hardware providing the massive parallel computation needed.

\begin{figure}[h]
\centering
\begin{tikzpicture}
\begin{axis}[
  width=10cm, height=6cm,
  xlabel={Amount of Training Data},
  ylabel={Model Performance},
  xmin=0, xmax=10, ymin=0, ymax=10,
  xtick=\empty, ytick=\empty,
  legend style={at={(0.97,0.25)}, anchor=east, font=\small, fill=white},
  axis lines=left,
  every axis plot/.append style={line width=1.4pt}]
\addplot[pOrange, domain=0:10, samples=200]{7*(1-exp(-0.8*x))};
\addlegendentry{Traditional ML}
\addplot[pBlue, dashed, domain=0:10, samples=200]{8*(1-exp(-0.5*x))};
\addlegendentry{Small neural network}
\addplot[pBlue, domain=0:10, samples=200]{9.5*(1-exp(-0.25*x))};
\addlegendentry{Large neural network}
\end{axis}
\end{tikzpicture}
\caption{Schematic performance versus training data size. Large neural networks
continue to benefit from additional data, whereas classical methods plateau.}
\label{fig:scaling}
\end{figure}

%----------------------------------------------------------
\section{Neural Network Model}

\subsection{The Single Neuron}

A single artificial neuron applies a \textbf{sigmoid activation function}:
\begin{equation}
a=g(z)=\frac{1}{1+e^{-z}},\qquad z=wx+b.
\label{eq:single-neuron}
\end{equation}
The term \emph{activation} is borrowed from neuroscience: just as a biological
neuron fires when stimulated strongly enough, an artificial neuron emits a
large output when its weighted input is large.

\subsection{Layered Architecture}

When multiple neurons are combined in a layered structure, the network can
capture complex non-linear relationships. A standard feedforward network has:
\begin{itemize}
\item \textbf{Input layer (Layer 0)}: the raw feature vector
  $\vec{x}\in\R^n$, also denoted $\va^{[0]}$.
\item \textbf{Hidden layers}: one or more intermediate layers. Each neuron in
  a hidden layer is a logistic-regression unit whose inputs are the activations
  of the previous layer.
\item \textbf{Output layer}: produces the final prediction.
\end{itemize}
The term \textit{hidden} means these layers are not directly observed in the
training data — we see only inputs $\vec{x}$ and labels $y$.

\subsection{Notation}

\begin{itemize}
\item Superscript $[l]$ denotes the \textbf{layer index}: $[1],[2],\ldots$
\item Subscript $j$ denotes the \textbf{neuron index} within a layer.
\item $a_j^{[l]}$: activation of neuron $j$ in layer $l$.
\item $\vw_j^{[l]}\in\R^{n_{l-1}}$, $b_j^{[l]}\in\R$: weight vector and bias.
\item $\va^{[l]}$: vector of all activations in layer $l$.
\item $\va^{[0]}=\vec{x}$ by definition.
\end{itemize}

\subsection{The General Activation Formula}

For any layer $l$ and unit $j$:
\begin{equation}
\boxed{a_j^{[l]}=g\!\Bigl(\vw_j^{[l]}\cdot\va^{[l-1]}+b_j^{[l]}\Bigr).}
\label{eq:activation}
\end{equation}
This single formula, applied repeatedly across layers and units, describes the
entire forward computation of a feedforward neural network.

\subsection{Layered Computation: Example}

Consider a network with architecture $n=4$ inputs, hidden layers of sizes 3
and 3, and a single output.

\textbf{Layer 1} (3 neurons, input $\va^{[0]}=\vec{x}$):
\begin{align*}
a_1^{[1]}&=g(\vw_1^{[1]}\cdot\vec{x}+b_1^{[1]}),\\
a_2^{[1]}&=g(\vw_2^{[1]}\cdot\vec{x}+b_2^{[1]}),\\
a_3^{[1]}&=g(\vw_3^{[1]}\cdot\vec{x}+b_3^{[1]}).
\end{align*}
Layer output: $\va^{[1]}=[a_1^{[1]},a_2^{[1]},a_3^{[1]}]^{\top}$.

\textbf{Layer 2} (output, 1 neuron):
\[a_1^{[2]}=g(\vw_1^{[2]}\cdot\va^{[1]}+b_1^{[2]})\in(0,1).\]

\textbf{Prediction}: $\hat{y}=1$ if $a_1^{[2]}\ge 0.5$, else $\hat{y}=0$.

\subsection{Deep Architectures and Hierarchical Features}

Networks with more than one hidden layer are called \textbf{deep networks},
and the field studying them is \textbf{deep learning}. In computer vision,
visualisation studies reveal a remarkable hierarchy of learned features:

\begin{center}
\begin{tabular}{lll}
\toprule
\textbf{Layer} & \textbf{Features detected} & \textbf{Description}\\
\midrule
1st hidden & Edges, oriented lines & Short segments at various angles\\
2nd hidden & Object parts & Eyes, corners, wheels\\
3rd hidden & Coarse shapes & Full face or car outlines\\
Output     & Category/identity & Final classification probability\\
\bottomrule
\end{tabular}
\end{center}

Crucially, \textbf{no human defines these features} — the network discovers
them automatically by minimising prediction error.

%----------------------------------------------------------
\section{TensorFlow Implementation}

\subsection{Data Representation}

TensorFlow processes data in \textbf{two-dimensional matrices} (tensors). Always
use 2D arrays when passing data to layers (shape: samples $\times$ features).

\subsection{The Dense Layer and Sequential API}

\begin{lstlisting}[caption={Sequential model for coffee roasting classification}]
import numpy as np
from tensorflow.keras.models import Sequential
from tensorflow.keras.layers import Dense

# Define architecture
model = Sequential([
    Dense(units=3, activation='sigmoid'),   # Hidden layer: 3 neurons
    Dense(units=1, activation='sigmoid')    # Output layer: 1 neuron
])
\end{lstlisting}

\subsection{The Training Workflow}

After defining architecture, training follows three steps:
\begin{enumerate}
\item \textbf{Compile} — specify loss and optimiser:
\begin{lstlisting}
model.compile(loss='binary_crossentropy', optimizer='adam')
\end{lstlisting}
\item \textbf{Train} — fit the model:
\begin{lstlisting}
model.fit(X_train, y_train, epochs=100)
\end{lstlisting}
\item \textbf{Predict}:
\begin{lstlisting}
predictions = model.predict(X_new)
\end{lstlisting}
\end{enumerate}

\subsection{Complete Example: Digit Recognition}

For $8\times 8$ pixel handwritten digit images ($64$ input features):
\begin{lstlisting}[caption={Neural network for handwritten digit recognition}]
model = Sequential([
    Dense(units=25, activation='sigmoid'),  # Layer 1: 25 units
    Dense(units=15, activation='sigmoid'),  # Layer 2: 15 units
    Dense(units=1,  activation='sigmoid')   # Output:   1 unit
])
\end{lstlisting}
The network compresses the representation from 64 dimensions down to a single
probability.

%----------------------------------------------------------
\section{Forward Propagation in NumPy}

Understanding how to implement forward propagation manually is essential for
debugging and appreciating the linear algebra underlying every deep learning
framework.

\begin{lstlisting}[caption={Reusable \texttt{dense()} layer function}]
import numpy as np

def sigmoid(z):
    return 1.0 / (1.0 + np.exp(-z))

def dense(a_in, W, b, activation=sigmoid):
    """
    a_in : (n_in,)       activations from previous layer
    W    : (n_in, units) weight matrix (each COLUMN = one neuron)
    b    : (units,)      bias vector
    """
    n_units = W.shape[1]
    a_out   = np.zeros(n_units)
    for j in range(n_units):
        w_j       = W[:, j]
        z_j       = np.dot(w_j, a_in) + b[j]
        a_out[j]  = activation(z_j)
    return a_out

def sequential(x, W1, b1, W2, b2):
    """Full forward pass through a 2-layer network."""
    a1  = dense(x,  W1, b1)
    a2  = dense(a1, W2, b2)
    return a2
\end{lstlisting}

\subsection{Vectorised Implementation}

\begin{lstlisting}[caption={Vectorised dense layer using \texttt{np.matmul}}]
def dense_vectorised(A_in, W, b, activation=sigmoid):
    """
    A_in : (1, n_in) -- 2D row vector input
    W    : (n_in, units)
    b    : (1, units) or broadcastable
    Returns A_out : (1, units)
    """
    Z     = np.matmul(A_in, W) + b   # (1, units)
    A_out = activation(Z)             # element-wise
    return A_out
\end{lstlisting}

For a batch of $m$ examples, $A_{\text{in}}$ has shape $(m\times n_{\text{in}})$
and one \texttt{matmul} computes all $m$ predictions simultaneously.

%----------------------------------------------------------
\section{Speculations on Artificial General Intelligence}

Current neural networks are \textbf{Artificial Narrow Intelligence (ANI)}:
trained for a specific task. \textbf{Artificial General Intelligence (AGI)}
would be capable of any intellectual task a human can perform and remains an
unsolved problem. Credible expert estimates for AGI range from ``within the
decade'' to ``never, as currently conceived'' — highlighting how poorly we
understand both the requirements for general intelligence and the trajectory of
AI progress.

%----------------------------------------------------------
\section*{Chapter Summary}
\addcontentsline{toc}{section}{Chapter Summary}

\begin{itemize}
\item A neural network stacks layers of artificial neurons, each computing
  $a_j^{[l]}=g(\vw_j^{[l]}\cdot\va^{[l-1]}+b_j^{[l]})$.
\item Deep networks automatically learn hierarchical features from data without
  any human-defined feature engineering.
\item TensorFlow's Sequential API defines, compiles, and trains a network in a
  few lines; the training loop internally uses backpropagation.
\item NumPy-level implementation reveals the matrix multiply at the core of
  every deep learning framework.
\item Vectorisation (batch matrix operations) is essential for efficiency.
\end{itemize}

% =========================================================
%  Chapter 5: Neural Network Training
% =========================================================
\chapter{Neural Network Training}

%----------------------------------------------------------
\section{A Three-Step Training Framework}

Whether fitting a two-parameter logistic regression or a million-parameter deep
network, the workflow decomposes into exactly three conceptual steps:

\begin{center}
\begin{tabular}{clll}
\toprule
\textbf{Step} & \textbf{Goal} & \textbf{Logistic Regression} &
\textbf{Neural Network (TF)}\\
\midrule
1 & Define $\vec{x}\to\hat{y}$ & $f=\sigma(\vec{w}\cdot\vec{x}+b)$ &
\texttt{Sequential([Dense(...)])}\\
2 & Measure error & Logistic loss & \texttt{model.compile(loss=...)}\\
3 & Minimise error & Gradient descent & \texttt{model.fit(X,y,epochs=...)}\\
\bottomrule
\end{tabular}
\end{center}

\subsection{Step 1 --- Architecture}

\begin{lstlisting}[caption={Binary digit classifier: 3-layer network}]
from tensorflow.keras import Sequential
from tensorflow.keras.layers import Dense

model = Sequential([
    Dense(units=25, activation='sigmoid'),  # Hidden layer 1
    Dense(units=15, activation='sigmoid'),  # Hidden layer 2
    Dense(units=1,  activation='sigmoid'),  # Output layer
])
\end{lstlisting}
TensorFlow initialises all weight matrices and bias vectors automatically.

\subsection{Step 2 --- Compile: Specify the Loss Function}

\subsubsection{Binary Cross-Entropy Loss}
For $(\vec{x},y)$ with $y\in\{0,1\}$:
\begin{equation}
L\bigl(f,y\bigr)=-y\log f-(1-y)\log(1-f).
\end{equation}
Cost over all $m$ examples:
\begin{equation}
J(\mathbf{W},\mathbf{B})=\frac{1}{m}\sum_{i=1}^m L\!\bigl(f(\vec{x}^{(i)}),y^{(i)}\bigr).
\end{equation}
\begin{lstlisting}
from tensorflow.keras.losses import BinaryCrossentropy
model.compile(loss=BinaryCrossentropy())
\end{lstlisting}

\subsubsection{Mean Squared Error (Regression)}
\begin{lstlisting}
from tensorflow.keras.losses import MeanSquaredError
model.compile(loss=MeanSquaredError())
\end{lstlisting}

\subsection{Step 3 --- Train: Backpropagation}

To minimise $J$, every parameter in every layer is updated iteratively:
\begin{align}
w_j^{[l]}&\leftarrow w_j^{[l]}-\alpha\,\pd{J}{w_j^{[l]}},\\
b_j^{[l]}&\leftarrow b_j^{[l]}-\alpha\,\pd{J}{b_j^{[l]}}.
\end{align}
Computing these partial derivatives across all layers is
\textbf{backpropagation} (Section~\ref{sec:backprop}).

\begin{definition}[Epoch]
One \textbf{epoch} is a single complete pass of gradient descent over the
entire training set.
\end{definition}
\begin{lstlisting}
model.fit(X, Y, epochs=100)
\end{lstlisting}

%----------------------------------------------------------
\section{Activation Functions}

\subsection{Why Not Just Use Sigmoid Everywhere?}

Using sigmoid in every hidden layer causes two problems:
\begin{enumerate}
\item \textbf{Range constraint.} Sigmoid outputs lie in $(0,1)$, but
  hidden-layer activations may represent quantities with unbounded range.
\item \textbf{Vanishing gradients.} At large $|z|$, $g'(z)\approx 0$, causing
  gradients to shrink exponentially during backpropagation through many layers.
\end{enumerate}

\subsection{Catalogue of Common Activation Functions}

\subsubsection{Sigmoid}
\begin{equation}
g(z)=\frac{1}{1+e^{-z}},\quad g(z)\in(0,1).
\end{equation}
Natural choice for binary classifier output.

\subsubsection{ReLU (Rectified Linear Unit)}
\begin{equation}
g(z)=\max(0,z),\quad g(z)\in[0,\infty).
\end{equation}
The \textbf{default for hidden layers} in modern deep learning. Advantages:
\begin{itemize}
\item \textbf{Computationally cheap}: a single comparison, no exponentiation.
\item \textbf{Non-vanishing gradients}: for $z>0$ the gradient is exactly $1$.
\end{itemize}

\subsubsection{Linear (Identity)}
\begin{equation}
g(z)=z,\quad g(z)\in(-\infty,+\infty).
\end{equation}
Appropriate \emph{only} in the output layer for regression (any-sign $y$).

\begin{figure}[h]
\centering
\begin{tikzpicture}
\begin{axis}[
  width=0.88\textwidth, height=5.8cm,
  xlabel={$z$}, ylabel={$g(z)$},
  xmin=-5, xmax=5, ymin=-1.2, ymax=3.5,
  xtick={-4,-2,0,2,4},
  legend pos=north west, legend style={font=\small, fill=white},
  grid=both, grid style={gray!15, thin}, axis lines=left,
  every axis plot/.append style={line width=1.4pt}]
\addplot[pBlue, domain=-5:5, samples=200]{1/(1+exp(-x))};
\addlegendentry{Sigmoid $\sigma(z)$}
\addplot[pRed, domain=-5:0, samples=2]{0};
\addplot[pRed, domain=0:5, samples=2, forget plot]{x};
\addlegendentry{ReLU $\max(0,z)$}
\addplot[pGreen, domain=-1.2:3.5, samples=2, dashed]{x};
\addlegendentry{Linear $z$}
\draw[pGray, dashed, thin] (axis cs:0,-1.2)--(axis cs:0,3.5);
\end{axis}
\end{tikzpicture}
\caption{Comparison of the three most common activation functions.}
\label{fig:activations}
\end{figure}

\subsection{Choosing the Right Activation}

\subsubsection{Output Layer --- Determined by Task}
\begin{center}
\begin{tabular}{llll}
\toprule
\textbf{Task} & \textbf{Target $y$} & \textbf{Activation} & \textbf{Range}\\
\midrule
Binary classification & $y\in\{0,1\}$ & Sigmoid  & $(0,1)$\\
Regression (any sign) & $y\in\R$       & Linear   & $(-\infty,+\infty)$\\
Regression ($\ge 0$)  & $y\ge 0$       & ReLU     & $[0,+\infty)$\\
\bottomrule
\end{tabular}
\end{center}

\subsubsection{Hidden Layers --- ReLU as Default}

\begin{lstlisting}[caption={Recommended: ReLU hidden layers with sigmoid output}]
model = Sequential([
    Dense(units=25, activation='relu'),
    Dense(units=15, activation='relu'),
    Dense(units=1,  activation='sigmoid'),
])
\end{lstlisting}

\subsection{The Necessity of Non-Linearity}

What happens if we use \emph{linear} activations everywhere?

\paragraph{Mathematical proof (one hidden unit).}
\begin{align*}
a^{[1]}&=w_1 x+b_1,\\
a^{[2]}&=w_2 a^{[1]}+b_2=w_2(w_1 x+b_1)+b_2=(w_2 w_1)x+(w_2 b_1+b_2)=Wx+B.
\end{align*}
Setting $W=w_2 w_1$ and $B=w_2 b_1+b_2$ gives a \emph{linear} model. The
hidden layer contributed nothing.

\textbf{Generalisation}: any depth of linear layers is equivalent to a single
linear layer. Adding more linear layers provides \emph{zero} additional
expressive power. Non-linear activations (ReLU) are what allow deep networks
to learn hierarchical, complex features.

\paragraph{Rule of thumb.}
\begin{itemize}
\item \textbf{Never} use linear activations in hidden layers.
\item Use \textbf{ReLU} for hidden layers.
\item Use \textbf{linear} in the output only for regression with $y\in\R$.
\end{itemize}

%----------------------------------------------------------
\section{Multiclass Classification}

\subsection{Definition and Motivation}

\begin{definition}[Multiclass classification]
A problem is \textbf{multiclass} when $y\in\{1,2,\ldots,N\}$ for $N>2$.
\end{definition}

Examples: handwritten digit recognition ($N=10$), medical diagnosis with
multiple disease types, part-of-speech tagging.

\subsection{Softmax Regression}

Logistic regression is a special case of \emph{softmax regression} with $N=2$.
For general $N$, each class $j$ has its own parameter vector $(\vec{w}_j,b_j)$.

\paragraph{Step 1: linear terms.}
\begin{equation}
z_j=\vec{w}_j\cdot\vec{x}+b_j,\quad j=1,\ldots,N.
\end{equation}

\paragraph{Step 2: softmax activation.}
\begin{equation}
\boxed{a_j=\frac{e^{z_j}}{\displaystyle\sum_{k=1}^N e^{z_k}}=P(y=j\mid\vec{x}).}
\label{eq:softmax}
\end{equation}

\paragraph{Normalisation property.}
\begin{equation}
\sum_{j=1}^N a_j=1.
\end{equation}

\begin{remark}[Coupling across units]
Unlike sigmoid or ReLU, softmax is \textbf{not element-wise}. Computing $a_j$
requires all $N$ values of $z$ simultaneously due to the shared denominator.
\end{remark}

\subsection{Cross-Entropy Loss for Softmax}

If the ground-truth label for example $i$ is $y^{(i)}=j$:
\begin{equation}
L=-\log(a_j)\quad\text{if }y=j.
\end{equation}
The total cost: $J=\frac{1}{m}\sum_{i=1}^m L(f(\vec{x}^{(i)}),y^{(i)})$.

\subsection{TensorFlow Implementation}

\begin{lstlisting}[caption={Softmax network for 10-class digit recognition}]
from tensorflow.keras.losses import SparseCategoricalCrossentropy

model = Sequential([
    Dense(units=25, activation='relu'),
    Dense(units=15, activation='relu'),
    Dense(units=10, activation='softmax'),   # Softmax output
])
model.compile(loss=SparseCategoricalCrossentropy())
model.fit(X, Y, epochs=100)
\end{lstlisting}

%----------------------------------------------------------
\section{Advanced Concepts}

\subsection{Numerical Stability: The \texttt{from\_logits} Pattern}

When TensorFlow evaluates softmax then log inside the cross-entropy loss, it
may suffer from floating-point overflow or underflow. The solution: let the
final layer output raw \textbf{logits} (the $z$ values) and pass
\lstinline{from_logits=True}. TensorFlow then combines the operations into a
single numerically stable formula.

\begin{definition}[Logit]
A \textbf{logit} is the raw, unscaled output $z$ of the final layer before any
activation.
\end{definition}

\begin{lstlisting}[caption={Numerically stable multiclass model (recommended in production)}]
model = Sequential([
    Dense(units=25, activation='relu'),
    Dense(units=15, activation='relu'),
    Dense(units=10, activation='linear'),   # Output logits, not probabilities
])
model.compile(loss=SparseCategoricalCrossentropy(from_logits=True))
model.fit(X, Y, epochs=100)

# Convert logits to probabilities at inference time:
import tensorflow as tf
logits = model.predict(X)
probs  = tf.nn.softmax(logits)
\end{lstlisting}

\subsection{Multi-Label Classification}

\begin{definition}[Multi-label classification]
The model must predict a \emph{vector} of binary labels; multiple labels can
be simultaneously active (non-mutually exclusive).
\end{definition}

\begin{example}
In autonomous driving, a single image may contain a car ($y_1=1$), no bus
($y_2=0$), and a pedestrian ($y_3=1$) simultaneously.
\end{example}

\begin{center}
\begin{tabular}{lll}
\toprule
\textbf{Feature} & \textbf{Multi-class} & \textbf{Multi-label}\\
\midrule
Output $y$           & Single integer   & Vector of $N$ binaries\\
Classes              & Mutually exclusive & Independent\\
Output activation    & Softmax          & Sigmoid (per output)\\
Loss                 & Sparse categorical CE & Binary CE\\
\bottomrule
\end{tabular}
\end{center}

\begin{lstlisting}[caption={Multi-label classification model}]
model = Sequential([
    Dense(units=25, activation='relu'),
    Dense(units=15, activation='relu'),
    Dense(units=3,  activation='sigmoid'),   # 3 independent binary outputs
])
model.compile(loss='binary_crossentropy')
\end{lstlisting}

\subsection{The Adam Optimiser}

Standard gradient descent uses a single fixed $\alpha$ for all parameters.

\begin{definition}[Adam]
\textbf{Adam} (Adaptive Moment Estimation) maintains a \emph{separate,
adaptive learning rate} $\alpha_j$ for every parameter $w_j$.
\end{definition}

Intuition:
\begin{itemize}
\item If $w_j$ consistently moves in the \emph{same direction}: increase $\alpha_j$.
\item If $w_j$ \emph{oscillates}: decrease $\alpha_j$.
\end{itemize}
Adam maintains two running statistics (moments) per parameter: a first moment
(exponentially weighted average of gradients) and a second moment (weighted
average of squared gradients). Their ratio yields a per-parameter step size.

\begin{lstlisting}[caption={Using Adam in TensorFlow}]
from tensorflow.keras.optimizers import Adam
model.compile(
    optimizer=Adam(learning_rate=1e-3),
    loss=SparseCategoricalCrossentropy(from_logits=True))
model.fit(X, Y, epochs=100)
\end{lstlisting}

%----------------------------------------------------------
\section{Backpropagation}
\label{sec:backprop}

\subsection{The Intuition Behind Derivatives}

Backpropagation computes the partial derivative of the cost $J$ with respect
to every parameter. These derivatives are used by the optimiser to update
parameters.

\begin{definition}[Derivative (informal)]
$\pd{J}{w}\approx k$ means: increasing $w$ by $\varepsilon$ causes $J$ to
change by approximately $k\varepsilon$.
\end{definition}

\paragraph{Example: $J(w)=w^2$.}
\begin{center}
\begin{tabular}{ccccc}
\toprule
$w$ & $J=w^2$ & $w+\varepsilon$ & $\Delta J$ & $\pd{J}{w}=2w$\\
\midrule
$3$  & $9$ & $3.001$ & $\approx 6\varepsilon$ & $6$\\
$2$  & $4$ & $2.001$ & $\approx 4\varepsilon$ & $4$\\
$-3$ & $9$ & $-2.999$ & $\approx -6\varepsilon$ & $-6$\\
\bottomrule
\end{tabular}
\end{center}

The derivative tells the optimiser: (1) which \emph{direction} to move, and
(2) how large a \emph{step} is appropriate.

\subsection{Computation Graphs and the Chain Rule}

\begin{definition}[Computation graph]
A \textbf{computation graph} is a DAG where each node represents an elementary
operation and each edge carries a value. Forward propagation traverses it
left-to-right; backpropagation right-to-left using the chain rule.
\end{definition}

\paragraph{Forward pass example.}
$a=wx+b$, $d=a-y$, $J=\tfrac{1}{2}d^2$.
With $x=-2$, $y=2$, $w=2$, $b=8$:
\[c=2(-2)=-4,\quad a=-4+8=4,\quad d=4-2=2,\quad J=2.\]

\paragraph{Backward pass (chain rule).}
\begin{align*}
\pd{J}{d}&=d=2,\quad
\pd{J}{a}=\pd{J}{d}\cdot 1=2,\\
\pd{J}{b}&=2,\quad
\pd{J}{c}=2,\quad
\pd{J}{w}=\pd{J}{c}\cdot x=2\cdot(-2)=-4.
\end{align*}

\subsection{Computational Efficiency}

\begin{center}
\begin{tabular}{lll}
\toprule
\textbf{Method} & \textbf{Complexity} & \textbf{Description}\\
\midrule
Naive numerical  & $O(N\times P)$ & Re-run forward pass for each parameter\\
Backpropagation  & $O(N+P)$       & Single backward pass, reuse partial derivatives\\
\bottomrule
\end{tabular}
\end{center}

For a network with one million parameters, the naive approach requires one
million forward passes per gradient step — physically intractable. Backpropagation
collapses this to a single backward pass of cost comparable to one forward pass.
Modern frameworks (TensorFlow, PyTorch) implement \textbf{automatic
differentiation}: they construct the computation graph dynamically and traverse
it in reverse, computing all gradients without any user-supplied formulas.

%----------------------------------------------------------
\section*{Chapter Summary}
\addcontentsline{toc}{section}{Chapter Summary}

\begin{center}
\begin{tabular}{ll}
\toprule
\textbf{Concept} & \textbf{Key takeaway}\\
\midrule
3-step training      & Specify $\to$ Compile $\to$ Fit\\
Activation choice    & ReLU for hidden; task-determined for output\\
Linear collapse      & All-linear networks $\equiv$ linear regression\\
Multiclass           & Softmax output with $N$ neurons\\
Numerical stability  & Use \texttt{from\_logits=True} in production\\
Multi-label          & $N$ independent sigmoid outputs\\
Adam                 & Per-parameter adaptive learning rates\\
Backpropagation      & $O(N+P)$ vs.\ $O(N\times P)$ naive; chain rule\\
Autodiff             & TF/PyTorch automate all gradient computation\\
\bottomrule
\end{tabular}
\end{center}

% =========================================================
%  Chapter 6: Advice for Applying Machine Learning
% =========================================================
\chapter{Advice for Applying Machine Learning}

\textit{This chapter provides a systematic framework for diagnosing and
improving machine learning systems. We cover practical debugging strategies,
the bias--variance trade-off, structured development workflows, data-centric
AI, and methods for handling skewed datasets.}

%----------------------------------------------------------
\section{Evaluating Model Performance}

\subsection{Unacceptable Prediction Errors: What Next?}

Suppose regularised linear regression produces unacceptably large errors.
You face a critical decision: \emph{what should you try next?} Choosing
wrongly can waste months of effort. Six commonly considered paths exist:

\begin{center}
\begin{tabular}{lp{9cm}}
\toprule
\textbf{Strategy} & \textbf{Description}\\
\midrule
Data collection      & Gather more training examples.\\
Feature reduction    & Use a smaller feature set.\\
Feature expansion    & Introduce additional informative features.\\
Polynomial complexity& Add polynomial terms for non-linear patterns.\\
Decrease $\lambda$   & Allow closer fitting of training data.\\
Increase $\lambda$   & Reduce overfitting by stronger regularisation.\\
\bottomrule
\end{tabular}
\end{center}

Without guidance, practitioners often choose randomly. \textbf{Machine
learning diagnostics} replace guessing with evidence-based reasoning.

\begin{definition}[Diagnostic]
A \textbf{diagnostic} is a formal test that provides insight into what is
or is not working within a learning algorithm, guiding targeted improvement.
\end{definition}

\subsection{The Train / Test Split}

The dataset is divided into two non-overlapping subsets:
\begin{enumerate}
\item \textbf{Training set} ($\approx70$--$80\%$): used to fit parameters $(\vec{w},b)$.
\item \textbf{Test set} ($\approx20$--$30\%$): used to evaluate on unseen data.
\end{enumerate}

\paragraph{Procedure for linear regression.}
\begin{enumerate}
\item Minimise the regularised cost on the training set.
\item Compute train and test errors \emph{without} regularisation:
\begin{align}
J_{\text{test}}&=\frac{1}{2m_{\text{test}}}\sum_{i=1}^{m_{\text{test}}}
\bigl(f_{\vec{w},b}(\vec{x}_{\text{test}}^{(i)})-y_{\text{test}}^{(i)}\bigr)^2,\\
J_{\text{train}}&=\frac{1}{2m_{\text{train}}}\sum_{i=1}^{m_{\text{train}}}
\bigl(f_{\vec{w},b}(\vec{x}_{\text{train}}^{(i)})-y_{\text{train}}^{(i)}\bigr)^2.
\end{align}
\end{enumerate}

\begin{remark}[Overfitting diagnostic rule]
If $J_{\text{train}}$ is low but $J_{\text{test}}$ is high, the model has
\textbf{failed to generalise} (overfitting).
\end{remark}

\subsection{The Three-Way Data Split for Model Selection}

A two-way split has a flaw: if we use the test set to select the best model
(e.g.\ the best polynomial degree $d$), $J_{\text{test}}$ is no longer an
unbiased estimate — it has been ``used up'' in the selection decision.

\textbf{Solution}: split data into three disjoint subsets.

\begin{center}
\begin{tabular}{lll}
\toprule
\textbf{Subset} & \textbf{Size} & \textbf{Purpose}\\
\midrule
Training set         & 60\% & Fit parameters $\vec{w}$, $b$\\
Cross-validation (CV) & 20\% & Select model / tune hyperparameters\\
Test set             & 20\% & Unbiased final performance estimate\\
\bottomrule
\end{tabular}
\end{center}

\paragraph{Correct procedure.}
\begin{enumerate}
\item Train each candidate model on the training set.
\item Evaluate all models on the CV set; choose the one with lowest $J_{\text{cv}}$.
\item Report $J_{\text{test}}$ of the \emph{single} chosen model.
\end{enumerate}

%----------------------------------------------------------
\section{Bias and Variance}

\subsection{Defining Bias and Variance}

\begin{definition}[High Bias]
A model has \textbf{high bias} when it is too simple to capture the underlying
patterns. Both $J_{\text{train}}$ and $J_{\text{cv}}$ are high, and close
to each other.
\end{definition}

\begin{definition}[High Variance]
A model has \textbf{high variance} when it fits training noise and fails to
generalise. $J_{\text{train}}$ is very low but
$J_{\text{cv}}\gg J_{\text{train}}$.
\end{definition}

\begin{center}
\begin{tabular}{llll}
\toprule
\textbf{Condition} & $J_{\text{train}}$ & $J_{\text{cv}}$ & \textbf{Relationship}\\
\midrule
High Bias     & High & High & $J_{\text{cv}}\approx J_{\text{train}}$\\
High Variance & Low  & High & $J_{\text{cv}}\gg J_{\text{train}}$\\
Just Right    & Low  & Low  & $J_{\text{cv}}\approx J_{\text{train}}$\\
Both High     & High & Much higher & Large gap, $J_{\text{train}}$ also high\\
\bottomrule
\end{tabular}
\end{center}

\begin{figure}[h]
\centering
\begin{tikzpicture}
\begin{axis}[
  width=10cm, height=6.5cm,
  xlabel={Polynomial degree $d$}, ylabel={Error},
  xmin=0.5, xmax=7.5, ymin=0, ymax=1.12,
  xtick={1,2,3,4,5,6,7}, ytick=\empty,
  legend pos=north east, legend style={font=\small, fill=white},
  axis lines=left,
  every axis plot/.append style={line width=1.4pt}]
\addplot[pBlue, smooth] coordinates
  {(1,.95)(2,.6)(3,.35)(4,.18)(5,.09)(6,.04)(7,.02)};
\addlegendentry{$J_{\text{train}}$}
\addplot[pOrange, dashed, smooth] coordinates
  {(1,1.0)(2,.55)(3,.3)(4,.36)(5,.56)(6,.79)(7,1.02)};
\addlegendentry{$J_{\text{cv}}$}
\draw[pGray, thin, dashed] (axis cs:3,0)--(axis cs:3,1.08);
\node[font=\small, pGray] at (axis cs:3,1.12) {Optimal $d$};
\node[pRed, font=\small] at (axis cs:1.4,.88) {High Bias};
\node[pRed, font=\small] at (axis cs:6.1,.90) {High Variance};
\end{axis}
\end{tikzpicture}
\caption{$J_{\text{train}}$ decreases monotonically with model complexity.
$J_{\text{cv}}$ follows a U-shaped curve with minimum at the optimal degree.}
\label{fig:complexity}
\end{figure}

\subsection{Regularisation and the Bias--Variance Trade-off}

\begin{itemize}
\item \textbf{Large $\lambda$}: forces $\vec{w}\approx\mathbf{0}$, model near
  constant $\Rightarrow$ \textbf{High Bias}.
\item \textbf{Small $\lambda$}: no constraint $\Rightarrow$ \textbf{High Variance}.
\item \textbf{Optimal $\lambda$}: balance between fitting data and simplicity.
\end{itemize}

\begin{figure}[h]
\centering
\begin{tikzpicture}
\begin{axis}[
  width=10cm, height=6.5cm,
  xlabel={$\lambda$}, ylabel={Error},
  xmin=-0.1, xmax=10.5, ymin=0, ymax=1.12,
  xtick={0,2,4,6,8,10}, ytick=\empty,
  legend pos=north east, legend style={font=\small, fill=white},
  axis lines=left,
  every axis plot/.append style={line width=1.4pt}]
\addplot[pBlue, smooth] coordinates
  {(0,.02)(1,.08)(2,.18)(3,.3)(5,.52)(7,.72)(10,.95)};
\addlegendentry{$J_{\text{train}}$}
\addplot[pOrange, dashed, smooth] coordinates
  {(0,1.0)(1,.55)(2,.30)(3,.29)(4,.33)(5,.46)(7,.66)(10,.96)};
\addlegendentry{$J_{\text{cv}}$}
\draw[pGray, thin, dashed] (axis cs:3,0)--(axis cs:3,1.08);
\node[font=\small, pGray] at (axis cs:3,1.12) {Optimal $\lambda$};
\node[pRed, font=\small] at (axis cs:0.7,.92) {High Variance};
\node[pRed, font=\small] at (axis cs:8.5,.78) {High Bias};
\end{axis}
\end{tikzpicture}
\caption{As $\lambda$ increases, $J_{\text{train}}$ rises and $J_{\text{cv}}$
forms a U-shape. The valley identifies the optimal regularisation strength.}
\label{fig:lambda}
\end{figure}

\subsection{Baseline Performance and the Gap Framework}

Comparing $J_{\text{train}}$ in isolation can be misleading — many tasks have
an irreducible noise floor.

\begin{definition}[Bias gap and variance gap]
\[
\text{Bias Gap}=J_{\text{train}}-\text{Baseline},\qquad
\text{Variance Gap}=J_{\text{cv}}-J_{\text{train}}.
\]
Large bias gap $\Rightarrow$ High Bias. Large variance gap $\Rightarrow$ High
Variance.
\end{definition}

\begin{example}[Speech recognition]
$J_{\text{train}}=10.8\%$, $J_{\text{cv}}=14.8\%$, human baseline $=10.6\%$.
\begin{itemize}
\item Bias Gap $=0.2\%$ (small).
\item Variance Gap $=4.0\%$ (large).
\item Conclusion: \textbf{High Variance} problem.
\end{itemize}
\end{example}

\subsection{Debugging Strategies}

\begin{center}
\begin{tabular}{ll}
\toprule
\textbf{Diagnosis} & \textbf{Recommended action}\\
\midrule
High Variance & Get more training examples\\
              & Reduce feature set\\
              & Increase $\lambda$\\
\midrule
High Bias     & Add more / better features\\
              & Add polynomial features\\
              & Decrease $\lambda$\\
\bottomrule
\end{tabular}
\end{center}

\subsection{Neural Networks and Bias--Variance}

A sufficiently large, well-regularised neural network can be made effectively
low-bias. The iterative workflow:
\begin{enumerate}
\item Check $J_{\text{train}}$ vs.\ baseline. If high $\Rightarrow$ High Bias:
  \textbf{use a bigger network} (more layers/units).
\item Check $J_{\text{cv}}$. If $J_{\text{cv}}\gg J_{\text{train}}$
  $\Rightarrow$ High Variance: \textbf{get more data}.
\item Repeat until both errors are acceptably low.
\end{enumerate}

\begin{lstlisting}[caption={$L_2$ regularisation in Keras}]
from tensorflow.keras.regularizers import L2
model = Sequential([
    Dense(units=25, activation='relu', kernel_regularizer=L2(0.01)),
    Dense(units=15, activation='relu', kernel_regularizer=L2(0.01)),
    Dense(units=1,  activation='sigmoid', kernel_regularizer=L2(0.01)),
])
\end{lstlisting}

%----------------------------------------------------------
\section{Machine Learning Development Process}

\subsection{The Iterative Development Loop}

Building a production ML model is rarely linear. It follows a continuous cycle:
\begin{enumerate}
\item \textbf{Choose architecture}: model family, features, hyperparameters.
\item \textbf{Train model}: implement and train.
\item \textbf{Run diagnostics}: bias--variance analysis and error analysis.
  Use findings to loop back.
\end{enumerate}

\subsection{Data Augmentation and Synthesis}

\begin{definition}[Data augmentation]
\textbf{Data augmentation} creates new training examples $(x',y)$ by applying
structure-preserving transformations to existing examples $(x,y)$.
\end{definition}

\begin{itemize}
\item \textbf{Computer vision}: rotation, cropping, mirroring, contrast
  changes, grid-based warping.
\item \textbf{Speech}: adding background noise, applying filters to simulate
  poor phone connections.
\end{itemize}

\textbf{Key rule}: Augmentation distortions must be \emph{representative of
real test-set noise}. Purely random pixel noise provides little benefit.

\begin{definition}[Data synthesis]
\textbf{Data synthesis} generates entirely new training examples from scratch
(e.g.\ rendering text in diverse fonts for OCR training).
\end{definition}

\subsection{Transfer Learning}

\begin{definition}[Transfer learning]
\textbf{Transfer learning} reuses a model pre-trained on a large related task
as the starting point for a new, data-scarce task.
\end{definition}

Standard two-step workflow:
\begin{enumerate}
\item \textbf{Supervised pre-training}: train a large network on a massive
  labelled dataset (e.g.\ 1 million ImageNet images).
\item \textbf{Fine-tuning}: replace the output layer for the target task. Then:
  \begin{itemize}
  \item \textbf{Frozen layers (Option 1)}: train only the new output layer.
    Best for very small new datasets ($\approx50$--$100$ examples).
  \item \textbf{Full fine-tuning (Option 2)}: retrain all layers using
    pre-trained weights as initialisation. Best for moderate-to-large
    new datasets.
  \end{itemize}
\end{enumerate}

Transfer learning works because neural networks learn a hierarchy of features:
early layers detect universal low-level patterns (edges, corners) that are
task-agnostic.

\subsection{Ethics and Fairness}

As ML systems impact billions of users, \textbf{fairness}, \textbf{bias
detection}, and \textbf{ethical development} are mandatory parts of the
lifecycle. A framework:
\begin{enumerate}[label=\Roman*.]
\item \textbf{Team diversity}: assemble diverse teams.
\item \textbf{Literature review}: consult industry guidelines and standards.
\item \textbf{Pre-deployment auditing}: measure performance across demographic
  subgroups.
\item \textbf{Mitigation and monitoring}: prepare rollback strategy; monitor
  for real-world harms post-deployment.
\end{enumerate}

%----------------------------------------------------------
\section{Skewed Datasets}

\subsection{The Problem with Accuracy on Skewed Data}

\begin{definition}[Skewed dataset]
A dataset is \textbf{skewed} when the ratio of positive to negative examples
is far from 50/50 (e.g.\ $99\%$ negative, $1\%$ positive).
\end{definition}

A classifier that \emph{always} predicts negative achieves $99\%$ accuracy
on such a dataset while providing zero useful predictions. This motivates
\textbf{Precision} and \textbf{Recall}.

\subsection{Confusion Matrix, Precision, and Recall}

\begin{center}
\begin{tabular}{llcc}
\toprule
& & \multicolumn{2}{c}{\textbf{Actual Class}}\\
\cmidrule{3-4}
& & Positive ($y=1$) & Negative ($y=0$)\\
\midrule
\multirow{2}{*}{\textbf{Predicted}} & Positive ($\hat{y}=1$) &
True Positive (TP) & False Positive (FP)\\
& Negative ($\hat{y}=0$) & False Negative (FN) & True Negative (TN)\\
\bottomrule
\end{tabular}
\end{center}

\begin{align}
\text{Precision}&=\frac{TP}{TP+FP},\qquad
\text{Recall}=\frac{TP}{TP+FN}.
\end{align}

A classifier that always predicts negative has $\text{Recall}=0$, correctly
exposing its uselessness.

\subsection{The Precision--Recall Trade-off}

Adjusting the decision threshold $\tau$ navigates the trade-off:
\begin{itemize}
\item Higher $\tau$ ($0.7$--$0.9$): higher precision, lower recall.
  Appropriate when false positives are costly.
\item Lower $\tau$ ($0.3$): higher recall, lower precision.
  Appropriate when false negatives are costly (e.g.\ missing a deadly disease).
\end{itemize}

\subsection{The F1 Score}

A single composite score combining precision and recall. The \emph{harmonic
mean} penalises extreme imbalance:
\begin{equation}
F_1=\frac{2PR}{P+R}.
\label{eq:f1}
\end{equation}
A high $F_1$ requires \emph{both} $P$ and $R$ to be high.

\begin{center}
\begin{tabular}{lcccc}
\toprule
\textbf{Algorithm} & $P$ & $R$ & \textbf{Arith.\ mean} & $F_1$\\
\midrule
Algorithm 1 & 0.50 & 0.40 & 0.450 & \textbf{0.444}\\
Algorithm 2 & 0.70 & 0.10 & 0.400 & 0.175\\
Algorithm 3 & 0.02 & 1.00 & 0.510 & 0.039\\
\bottomrule
\end{tabular}
\end{center}

Algorithm~3 achieves the highest arithmetic mean ($0.51$) but its F1 of $0.039$
correctly reveals it as a poor classifier (predicts positive for everything).

%----------------------------------------------------------
\section*{Chapter Summary}
\addcontentsline{toc}{section}{Chapter Summary}

\begin{itemize}
\item \textbf{Diagnostics} replace guessing with evidence. The three-way
  split (train/CV/test) enables unbiased model selection.
\item \textbf{Bias} (underfitting) and \textbf{variance} (overfitting) are the
  two primary failure modes, diagnosed by comparing $J_{\text{train}}$ vs.\
  $J_{\text{cv}}$ against a baseline.
\item The iterative development loop (architecture $\to$ train $\to$ diagnose)
  is the standard workflow. Large, well-regularised networks tend to be
  low-bias; more data cures high variance.
\item Data augmentation, synthesis, and transfer learning address data scarcity.
\item \textbf{Skewed datasets} require Precision, Recall, and $F_1$ score
  rather than raw accuracy.
\item Ethics and fairness must be embedded throughout the development lifecycle.
\end{itemize}

% =========================================================
%  Chapter 7: Decision Trees
% =========================================================
\chapter{Decision Trees}

Decision trees are among the most powerful and interpretable machine learning
algorithms in widespread use. Unlike linear models, they naturally capture
complex non-linear relationships without requiring feature scaling or
distributional assumptions, and they serve as the building block for advanced
ensemble methods --- Random Forests and XGBoost --- that consistently rank
among the top-performing algorithms on structured tabular data.

\begin{remark}
Decision trees are particularly effective for \textbf{tabular (structured)
data}. For unstructured data (images, audio, raw text), neural networks are
generally preferred.
\end{remark}

%----------------------------------------------------------
\section{Decision Tree Basics}

\subsection{Dataset Representation}

Consider classifying animals as cats or not, given three observable features:
\begin{itemize}
\item $x_1$: Ear Shape --- Pointy or Floppy
\item $x_2$: Face Shape --- Round or Not Round
\item $x_3$: Whiskers --- Present or Absent
\item $y\in\{0,1\}$: $1=$Cat, $0=$Not Cat
\end{itemize}

Decision trees can handle both categorical and continuous features. For
categorical features with more than two values, one-hot encoding
(Section~\ref{sec:ohe}) converts each category into a binary indicator.

\subsection{Anatomy of a Decision Tree}

A decision tree is a hierarchical, acyclic graph encoding a sequence of
if--then--else rules.

\begin{description}
\item[Root Node] The topmost node. Receives the full training dataset.
\item[Internal Nodes] Test a specific feature and route examples to child
  branches based on the outcome.
\item[Leaf Nodes] Terminal nodes storing a class label (classification) or
  numeric value (regression).
\end{description}

The \textbf{depth} of a node is the number of edges from the root. The root
is at depth~$0$.

\subsection{Inference}

To classify a new example:
\begin{enumerate}
\item Start at the root node.
\item Evaluate the feature tested at the current node.
\item Follow the matching branch.
\item Repeat at each subsequent node.
\item Return the prediction stored in the reached leaf.
\end{enumerate}
Inference is $O(d)$ in tree depth $d$ --- very efficient.

%----------------------------------------------------------
\section{Decision Tree Learning}

\subsection{The Recursive Splitting Algorithm}

Tree construction proceeds top-down, greedy, and recursively:
\begin{enumerate}
\item Place the entire training set at the root.
\item Select the feature that maximises Information Gain.
\item Split examples into subsets by feature value.
\item Recursively repeat on each child subset.
\item Apply stopping criteria to decide when a node becomes a leaf.
\end{enumerate}

\subsection{Entropy: Measuring Impurity}
\label{sec:entropy}

At any node, the training subset may mix multiple classes. The degree of mixing
is called \textbf{impurity}. \textbf{Entropy} is the standard measure:

\begin{equation}
H(p_1)=-p_1\log_2 p_1-(1-p_1)\log_2(1-p_1),
\label{eq:entropy}
\end{equation}
where $p_1$ is the positive-class fraction. Convention: $0\log_2 0\equiv 0$.

Properties:
\begin{itemize}
\item $H(p_1)=0$ when $p_1\in\{0,1\}$ (pure node).
\item $H(p_1)=1$ when $p_1=0.5$ (maximum uncertainty).
\item $H$ is concave and symmetric, achieving its maximum at $p_1=0.5$.
\end{itemize}

\begin{figure}[h]
\centering
\begin{tikzpicture}
\begin{axis}[
  width=0.62\textwidth, height=5.5cm,
  xlabel={Positive-class fraction $p_1$}, ylabel={Entropy $H(p_1)$ (bits)},
  xmin=0, xmax=1, ymin=0, ymax=1.12,
  xtick={0,0.25,0.5,0.75,1}, ytick={0,0.25,0.5,0.75,1},
  grid=both, grid style={gray!15, thin},
  tick label style={font=\small}, label style={font=\small},
  axis lines=left,
  every axis plot/.append style={line width=1.5pt}]
\addplot[pBlue, domain=0.001:0.999, samples=200]
  {-x*log2(x)-(1-x)*log2(1-x)};
\addplot[only marks, mark=*, color=pRed, mark size=3.5pt]
  coordinates{(0.5,1.0)};
\node[above right, font=\small, pRed] at (axis cs:0.52,1.0) {max at $p_1=0.5$};
\end{axis}
\end{tikzpicture}
\caption{Binary entropy $H(p_1)$. It peaks at 1 bit (maximum uncertainty) when
the two classes are equally likely, and vanishes at pure nodes.}
\label{fig:entropy-curve}
\end{figure}

\begin{center}
\begin{tabular}{lccr}
\toprule
\textbf{Composition} & $p_1$ & $H(p_1)$ & \textbf{Impurity}\\
\midrule
6 cats, 0 dogs & 1.000 & 0.000 & None (pure)\\
5 cats, 1 dog  & 0.833 & 0.650 & Low\\
4 cats, 2 dogs & 0.667 & 0.918 & Moderate\\
3 cats, 3 dogs & 0.500 & 1.000 & Maximum\\
0 cats, 6 dogs & 0.000 & 0.000 & None (pure)\\
\bottomrule
\end{tabular}
\end{center}

\subsection{Information Gain}
\label{sec:infogain}

Let $n$ be the number of examples at the current node, $n_L$ and $n_R$ those
reaching the left and right children. The \textbf{Information Gain} is:
\begin{equation}
\IG(f)=H(p_1)-\Bigl[w^L H(p_1^L)+w^R H(p_1^R)\Bigr],\quad
w^L=\frac{n_L}{n},\;w^R=\frac{n_R}{n}.
\label{eq:ig}
\end{equation}
IG is always non-negative (by convexity of entropy) and equals zero when the
split provides no useful information.

\begin{example}
Root: 10 examples, 5 cats, 5 dogs $\Rightarrow H(p_1)=1$.

\textbf{Ear Shape:} Left (5 ex., 4 cats): $H=0.72$. Right (5 ex., 1 cat):
$H=0.72$. $\IG=1-[0.5(0.72)+0.5(0.72)]=\mathbf{0.28}$.

\textbf{Face Shape:} $\IG\approx\mathbf{0.03}$.

\textbf{Whiskers:} $\IG\approx\mathbf{0.12}$.

\textbf{Decision}: select \textit{Ear Shape} (largest $\IG$).
\end{example}

\subsection{Stopping Criteria}
\label{sec:stopping}

\begin{enumerate}
\item \textbf{Perfect purity}: stop when $H=0$.
\item \textbf{Maximum depth}: stop at depth $d_{\max}$.
\item \textbf{Minimum IG}: stop if best $\IG<\epsilon_{\IG}$.
\item \textbf{Minimum node size}: stop if $n<n_{\min}$.
\end{enumerate}

\subsection{One-Hot Encoding}
\label{sec:ohe}

\begin{definition}[One-Hot Encoding]
A categorical feature $x$ with $k$ values is replaced by $k$ binary features
$b_1,\ldots,b_k$, where $b_j=1$ if $x$ equals the $j$-th category.
\end{definition}

\begin{example}
Ear Shape $\in\{\text{Pointy, Floppy, Oval}\}$:
\begin{center}
\begin{tabular}{lccc}
\toprule
Shape & Is Pointy? & Is Floppy? & Is Oval?\\
\midrule
Pointy & 1 & 0 & 0\\
Floppy & 0 & 1 & 0\\
Oval   & 0 & 0 & 1\\
\bottomrule
\end{tabular}
\end{center}
\end{example}

After OHE, every feature is binary and the standard IG-based splitting applies.

\subsection{Regression Trees}

A \textbf{regression tree} predicts a continuous target. A regression leaf
stores the \textbf{mean} of the target values of all training examples it
receives:
\[
\hat{y}_{\text{leaf}}=\frac{1}{|\mathcal{S}_{\text{leaf}}|}
\sum_{i\in\mathcal{S}_{\text{leaf}}}y_i.
\]

The splitting criterion is \textbf{variance reduction}:
\begin{equation}
\Delta\sigma^2=\sigma^2_{\text{root}}
-\Bigl[\frac{n_L}{n}\,\sigma^2_L+\frac{n_R}{n}\,\sigma^2_R\Bigr].
\label{eq:var-reduction}
\end{equation}

\begin{center}
\begin{tabular}{lll}
\toprule
\textbf{Aspect} & \textbf{Classification tree} & \textbf{Regression tree}\\
\midrule
Target $y$       & Discrete label   & Continuous number\\
Leaf prediction  & Majority class   & Mean of examples\\
Split criterion  & Information Gain & Variance Reduction\\
\bottomrule
\end{tabular}
\end{center}

%----------------------------------------------------------
\section{Tree Ensembles}

\subsection{Motivation: Instability of Single Trees}

A single decision tree has \textbf{high variance}: a small perturbation in the
training data can lead to a completely different tree structure. The solution is
to combine many diverse trees.

\subsection{Bootstrapping and Bagging}

\textbf{Bootstrapping}: draw $m$ examples with replacement from a training set
of size $m$. The resulting sample contains some examples multiple times and
omits others. On average, each bootstrap sample includes $\approx63.2\%$ of the
unique examples.

\textbf{Bagging} (Bootstrap Aggregating): train an independent decision tree on
each of $B$ bootstrap samples, then aggregate predictions by majority vote
(classification) or averaging (regression).

\textbf{Hyperparameter $B$}: typical values $B\in[64,128]$. Increasing $B$
never hurts accuracy but yields diminishing returns beyond $\approx128$.

\subsection{Random Forests}
\label{sec:rf}

In a bagged ensemble all trees share the same full feature set; if one feature
is a strong predictor, nearly every tree will use it as the root split, making
trees highly correlated. \textbf{Random Forest} adds a second source of
randomness: at each node, only a random subset of $k<n$ features is evaluated:

\begin{equation}
k=\lfloor\sqrt{n}\rfloor\quad\text{(standard for classification)}.
\label{eq:rf-k}
\end{equation}

This \emph{decorrelates} the trees, amplifying the variance-reduction benefit
of averaging.

\subsection{Boosting and XGBoost}

\textbf{Bagging} trains trees \emph{in parallel} on independent subsamples.
\textbf{Boosting} trains trees \emph{sequentially}: each new tree focuses on
examples the current ensemble handles poorly. Key insight: concentrate effort
on weaknesses, not strengths (like deliberate practice).

\textbf{XGBoost} (eXtreme Gradient Boosting) frames boosting as gradient
descent in function space. Key engineering features:
\begin{itemize}
\item Built-in $L_1$ and $L_2$ regularisation on leaf weights.
\item Efficient $O(n\log n)$ split-finding via quantile sketches.
\item CPU cache-aware computation.
\item Native handling of missing values.
\end{itemize}

\begin{lstlisting}[caption={XGBoost for classification and regression}]
from xgboost import XGBClassifier, XGBRegressor

# Classification
clf = XGBClassifier(n_estimators=100, max_depth=6,
                    learning_rate=0.1, random_state=42)
clf.fit(X_train, y_train)
y_pred = clf.predict(X_test)

# Regression
reg = XGBRegressor(n_estimators=100, max_depth=6,
                   learning_rate=0.1, random_state=42)
reg.fit(X_train, y_train)
\end{lstlisting}

\subsection{Decision Trees vs.\ Neural Networks}

\begin{center}
\begin{tabular}{lp{5cm}p{5cm}}
\toprule
\textbf{Criterion} & \textbf{Decision Trees/Ensembles} & \textbf{Neural Networks}\\
\midrule
Best data type    & Tabular/structured         & Unstructured, mixed\\
Training speed    & Fast                        & Slow (large models)\\
Interpretability  & High (single trees)         & Low (black box)\\
Transfer learning & Not available               & Yes\\
Feature scaling   & Not required               & Recommended\\
Recommended tool  & XGBoost, Random Forest     & PyTorch, TensorFlow\\
\bottomrule
\end{tabular}
\end{center}

\textbf{Practical guidance}:
\begin{itemize}
\item Use XGBoost/Random Forest for tabular data, fast iteration, or when
  interpretability is required.
\item Use neural networks for images, text, audio, video, or mixed modalities.
\end{itemize}

%----------------------------------------------------------
\section*{Chapter Summary}
\addcontentsline{toc}{section}{Chapter Summary}

\textbf{Decision Tree Basics}: hierarchical if--then--else structure traversed
in $O(d)$ time; handles categorical and continuous features.

\textbf{Learning algorithm}: recursive top-down splitting that maximises
Information Gain $\IG(f)=H(p_1)-[w^L H(p_1^L)+w^R H(p_1^R)]$ at each node.
Regression trees replace entropy with variance reduction.

\textbf{Tree Ensembles}: bagging (bootstrap + average) reduces variance.
Random Forests add feature randomisation to decorrelate trees. XGBoost uses
sequential boosting with gradient descent in function space, achieving
state-of-the-art performance on structured data.

% =========================================================
%  Chapter 8: Unsupervised Learning
% =========================================================
\chapter{Unsupervised Learning}

%----------------------------------------------------------
\section{Clustering}

\subsection{Overview and Motivation}

\textbf{Clustering} is one of the foundational unsupervised learning
algorithms. Its goal is to automatically discover patterns, groupings, or
structures within a dataset --- without any human-provided labels.

\begin{definition}[Cluster]
A \textbf{cluster} is a subset of training examples such that points within
the subset are more similar to one another than to points outside it.
\end{definition}

Real-world applications:
\begin{enumerate}
\item \textbf{News aggregation}: grouping articles by topic.
\item \textbf{Market segmentation}: identifying distinct customer groups.
\item \textbf{Genomics}: grouping individuals by genetic expression data.
\item \textbf{Astronomy}: grouping celestial bodies into coherent structures.
\item \textbf{Image compression}: representing an image using $K$ colours.
\end{enumerate}

\begin{center}
\begin{tabular}{lp{5.2cm}p{5.2cm}}
\toprule
\textbf{Feature} & \textbf{Supervised learning} & \textbf{Clustering}\\
\midrule
Data format    & $(x,y)$ pairs with labels  & $x$ only\\
Goal           & Learn $f$ s.t.\ $f(x)\approx y$ & Discover groups\\
Error signal   & Labelled ground truth      & Cost / distortion function\\
\bottomrule
\end{tabular}
\end{center}

\subsection{The K-Means Algorithm}

K-means partitions $m$ examples into $K$ clusters by alternating two steps.

\paragraph{Initialisation.} Randomly select $K$ distinct training examples and
place initial centroids $\mu_1,\ldots,\mu_K$ at those locations.

\paragraph{Step 1 --- Cluster assignment.}
\[c^{(i)}:=\underset{k}{\argmin}\;\|x^{(i)}-\mu_k\|^2.\]

\paragraph{Step 2 --- Centroid update.}
\[\mu_k:=\frac{1}{|\mathcal{C}_k|}\sum_{i:\,c^{(i)}=k}x^{(i)}.\]

If $\mathcal{C}_k=\emptyset$, eliminate that cluster (reduce $K$ by one).

\begin{figure}[h]
\centering
\begin{subfigure}[b]{0.30\textwidth}
\centering
\begin{tikzpicture}[scale=0.6]
\draw[->](-0.3,0)--(5.8,0)node[right]{$x_1$};
\draw[->](0,-0.3)--(0,5.5)node[above]{$x_2$};
\foreach \x/\y in {0.6/1.2,1.0/0.8,1.4/1.5,0.8/2.0,1.2/0.5,
    2.8/3.8,3.2/4.2,2.5/3.5,3.5/3.8,3.0/4.5,
    4.2/1.0,4.5/1.8,4.0/0.8,4.8/1.4,4.3/2.2}{
  \fill[black!30](\x,\y)circle(2.5pt);}
\fill[cA](1.5,3.5)node[above right,font=\tiny]{$\mu_1$}circle(5pt);
\fill[cB](3.0,1.5)node[above right,font=\tiny]{$\mu_2$}circle(5pt);
\fill[cC](4.0,3.5)node[above right,font=\tiny]{$\mu_3$}circle(5pt);
\end{tikzpicture}
\caption{Initialisation}
\end{subfigure}
\hfill
\begin{subfigure}[b]{0.30\textwidth}
\centering
\begin{tikzpicture}[scale=0.6]
\draw[->](-0.3,0)--(5.8,0)node[right]{$x_1$};
\draw[->](0,-0.3)--(0,5.5)node[above]{$x_2$};
\foreach \x/\y in {0.6/1.2,1.0/0.8,1.4/1.5,0.8/2.0,1.2/0.5}
  {\fill[cA!80](\x,\y)circle(2.5pt);}
\foreach \x/\y in {4.2/1.0,4.5/1.8,4.0/0.8,4.8/1.4,4.3/2.2}
  {\fill[cB!80](\x,\y)circle(2.5pt);}
\foreach \x/\y in {2.8/3.8,3.2/4.2,2.5/3.5,3.5/3.8,3.0/4.5}
  {\fill[cC!80](\x,\y)circle(2.5pt);}
\fill[cA](1.5,3.5)circle(5pt);
\fill[cB](3.0,1.5)circle(5pt);
\fill[cC](4.0,3.5)circle(5pt);
\end{tikzpicture}
\caption{After assignment}
\end{subfigure}
\hfill
\begin{subfigure}[b]{0.30\textwidth}
\centering
\begin{tikzpicture}[scale=0.6]
\draw[->](-0.3,0)--(5.8,0)node[right]{$x_1$};
\draw[->](0,-0.3)--(0,5.5)node[above]{$x_2$};
\foreach \x/\y in {0.6/1.2,1.0/0.8,1.4/1.5,0.8/2.0,1.2/0.5}
  {\fill[cA!80](\x,\y)circle(2.5pt);}
\foreach \x/\y in {4.2/1.0,4.5/1.8,4.0/0.8,4.8/1.4,4.3/2.2}
  {\fill[cB!80](\x,\y)circle(2.5pt);}
\foreach \x/\y in {2.8/3.8,3.2/4.2,2.5/3.5,3.5/3.8,3.0/4.5}
  {\fill[cC!80](\x,\y)circle(2.5pt);}
\fill[cD](1.0,1.2)node[above right,font=\tiny]{$\mu_1$}circle(5pt);
\fill[cD](4.36,1.44)node[above right,font=\tiny]{$\mu_2$}circle(5pt);
\fill[cD](3.0,3.96)node[above right,font=\tiny]{$\mu_3$}circle(5pt);
\end{tikzpicture}
\caption{Converged}
\end{subfigure}
\caption{K-means with $K=3$. Coloured points are cluster assignments; orange
diamonds are centroids.}
\label{fig:kmeans}
\end{figure}

\subsection{The K-Means Cost Function}

K-means minimises the \textbf{distortion function}:
\[J=\frac{1}{m}\sum_{i=1}^m\|x^{(i)}-\mu_{c^{(i)}}\|^2.\]

Each step provably decreases $J$ or leaves it unchanged, so convergence is
guaranteed. An increasing $J$ signals a code bug.

\paragraph{Multiple initialisations.} Run K-means $R=50$--$1000$ times with
different random initialisations; keep the solution with the lowest $J$.
This is especially important for small $K$.

\subsection{Choosing $K$}

\paragraph{Elbow method.} Plot $J$ vs.\ $K$; look for a sharp bend. Often
unreliable because many real-world curves lack a clear elbow.

\paragraph{Downstream purpose.} Choose $K$ based on the end objective
(e.g.\ S/M/L for T-shirt sizing, or quality--compression trade-off for image
compression).

%----------------------------------------------------------
\section{Anomaly Detection}

\subsection{Introduction and Types of Anomalies}

\textbf{Anomaly detection} identifies data points that deviate significantly
from normal behaviour.

\begin{definition}[Anomaly]
$x_{\text{test}}$ is anomalous if $p(x_{\text{test}})<\epsilon$
under a model of normal data.
\end{definition}

Three canonical types:
\begin{itemize}
\item \textbf{Point anomalies}: a single point far from the distribution
  (e.g.\ a credit card transaction of \$15{,}000 for a customer spending
  \$500/month).
\item \textbf{Contextual anomalies}: normal in absolute terms but abnormal
  in context (e.g.\ $85^\circ$F in January in New York).
\item \textbf{Collective anomalies}: a group of points individually normal but
  anomalous together (e.g.\ thousands of packets from one IP in seconds,
  suggesting a DoS attack).
\end{itemize}

Applications: fraud detection, cybersecurity, manufacturing quality control,
healthcare monitoring, data centre monitoring.

\subsection{The Gaussian Density Model}

The dominant framework: fit a Gaussian to each feature, then flag examples
with very low joint probability.

If $X\sim\mathcal{N}(\mu,\sigma^2)$, its PDF is:
\[p(x;\mu,\sigma^2)=\frac{1}{\sqrt{2\pi}\,\sigma}
\exp\!\Bigl(-\frac{(x-\mu)^2}{2\sigma^2}\Bigr).\]

Maximum likelihood estimates from training data:
\[\hat{\mu}=\frac{1}{m}\sum_{i=1}^m x^{(i)},\qquad
\hat{\sigma}^2=\frac{1}{m}\sum_{i=1}^m(x^{(i)}-\hat{\mu})^2.\]

\subsection{Multi-Dimensional Algorithm}

Under the independence assumption, the joint density factors as:
\[p(\vec{x})=\prod_{j=1}^n p(x_j;\mu_j,\sigma_j^2).\]

\paragraph{Algorithm.}
\begin{enumerate}[label=\textbf{Step \arabic*:}]
\item Choose $n$ informative features.
\item Fit $\mu_j$ and $\sigma_j^2$ for each feature $j$ on the training set.
\item For a new example: compute $p(\vec{x})$; flag as anomaly if
  $p(\vec{x})<\epsilon$.
\end{enumerate}

\subsection{Evaluation on Skewed Data}

Use precision, recall, and $F_1$ (not accuracy). Choose $\epsilon$ to maximise
$F_1$ on the cross-validation set.

\subsection{Anomaly Detection vs.\ Supervised Learning}

\begin{center}
\begin{tabular}{lp{5.2cm}p{5.2cm}}
\toprule
\textbf{Feature} & \textbf{Anomaly Detection} & \textbf{Supervised Learning}\\
\midrule
Positive examples & Very few (0--20) & Many\\
Anomaly types & Diverse, unpredictable & Recurring, resembles training data\\
Use case & Fraud, novel defects & Spam, known defects\\
\bottomrule
\end{tabular}
\end{center}

\subsection{Feature Engineering Tips}

\begin{itemize}
\item Apply $\log(x)$ or $x^p$ transforms if a feature is heavily skewed.
\item Create ratio features (e.g.\ CPU load / network traffic) to capture
  anomalous relationships between features.
\item Perform error analysis: inspect missed anomalies to design new features.
\end{itemize}

%----------------------------------------------------------
\section*{Chapter Summary}
\addcontentsline{toc}{section}{Chapter Summary}

\paragraph{Clustering.} K-means alternates between assignment and centroid
update, minimising the distortion $J=\tfrac{1}{m}\sum\|x^{(i)}-\mu_{c^{(i)}}\|^2$.
Run multiple random initialisations. Choose $K$ based on downstream purpose.

\paragraph{Anomaly detection.} Fits a Gaussian density model to normal data
and flags examples with $p(\vec{x})<\epsilon$. Evaluation uses $F_1$; choose
$\epsilon$ on the CV set. Use anomaly detection when anomaly types are diverse
and few positive examples exist; use supervised learning when they are
recurring and plentiful.

% =========================================================
%  Chapter 9: Recommender Systems
% =========================================================
\chapter{Recommender Systems}

Recommender systems are among the most commercially successful applications of
machine learning. They power suggestions on streaming platforms such as Netflix
and Spotify, product recommendations on Amazon, and content feeds of social
media platforms. At their core, they learn from past interactions between users
and items to predict which items a user is most likely to enjoy next.

Formally: $n$ users, $m$ items, rating matrix $\mathbf{R}\in\R^{n\times m}$
where $r_{ui}$ is user $u$'s rating of item $i$. Because most users rate only
a tiny fraction of items, this matrix is extremely \emph{sparse}. The task is
to infer the missing entries.

%----------------------------------------------------------
\section{Collaborative Filtering}
\label{sec:cf}

\subsection{Core Idea}

Collaborative filtering's fundamental assumption: \emph{users who agreed in the
past tend to agree in the future}. No item content information is required.

\subsection{Memory-Based CF}

\subsubsection{User-Based CF}

Find users similar to the target user; use their ratings to predict:
\begin{equation}
\hat{r}_{ui}=\bar{r}_u+\frac{\sum_{v\in\mathcal{N}(u)}\mathrm{sim}(u,v)(r_{vi}-\bar{r}_v)}
{\sum_{v\in\mathcal{N}(u)}|\mathrm{sim}(u,v)|}.
\end{equation}

\subsubsection{Similarity Measures}

\paragraph{Cosine similarity.}
\[\mathrm{sim}_{\cos}(\mathbf{a},\mathbf{b})=\frac{\mathbf{a}\cdot\mathbf{b}}{\|\mathbf{a}\|_2\|\mathbf{b}\|_2}.\]

\paragraph{Pearson correlation.}
\[\mathrm{sim}_{\text{P}}(u,v)=\frac{\sum_{i\in\mathcal{I}_{uv}}(r_{ui}-\bar{r}_u)(r_{vi}-\bar{r}_v)}
{\sqrt{\sum_i(r_{ui}-\bar{r}_u)^2}\,\sqrt{\sum_i(r_{vi}-\bar{r}_v)^2}}.\]

\subsection{Model-Based CF: Matrix Factorisation}

Approximate the rating matrix as a product of two low-rank factor matrices:
\begin{equation}
\mathbf{R}\approx\mathbf{U}\mathbf{V}^{\!\top},
\label{eq:mf}
\end{equation}
$\mathbf{U}\in\R^{n\times k}$ (user factors), $\mathbf{V}\in\R^{m\times k}$
(item factors), $k\ll\min(n,m)$. Predicted rating:
\[\hat{r}_{ui}=\mathbf{u}_u\cdot\mathbf{v}_i.\]

\paragraph{Learning.} Minimise regularised MSE over observed ratings:
\begin{equation}
J(\mathbf{U},\mathbf{V})=\frac{1}{|\Omega|}\sum_{(u,i)\in\Omega}
(r_{ui}-\mathbf{u}_u\cdot\mathbf{v}_i)^2
+\lambda\!\Bigl(\sum_u\|\mathbf{u}_u\|^2+\sum_i\|\mathbf{v}_i\|^2\Bigr).
\end{equation}

SGD updates for observation $(u,i)$:
\begin{align}
e_{ui}&=r_{ui}-\mathbf{u}_u\cdot\mathbf{v}_i,\\
\mathbf{u}_u&\leftarrow\mathbf{u}_u+\alpha(e_{ui}\mathbf{v}_i-\lambda\mathbf{u}_u),\\
\mathbf{v}_i&\leftarrow\mathbf{v}_i+\alpha(e_{ui}\mathbf{u}_u-\lambda\mathbf{v}_i).
\end{align}

\paragraph{Bias terms.} An augmented model:
\[\hat{r}_{ui}=\mu+b_u+b_i+\mathbf{u}_u\cdot\mathbf{v}_i,\]
where $\mu$ is the global mean, $b_u$ the user bias, $b_i$ the item bias.

\subsection{Limitations of Collaborative Filtering}

\begin{itemize}
\item \textbf{Cold-start problem}: new users or items with no ratings.
\item \textbf{Popularity bias}: over-recommends popular items.
\item \textbf{Filter bubbles}: reinforces existing preferences.
\end{itemize}

%----------------------------------------------------------
\section{Implementation Details}

\subsection{Cost Function}

Let $n_u$ users, $n_m$ movies. Indicator $r^{(i,j)}=1$ if user $j$ rated
movie $i$.

\begin{equation}
J=\frac{1}{2}\sum_{(i,j):\,r^{(i,j)}=1}
\!\!\bigl(\vec{w}^{(j)}\cdot\vec{x}^{(i)}+b^{(j)}-y^{(i,j)}\bigr)^2
+\frac{\lambda}{2}\sum_j\sum_k(w_k^{(j)})^2
+\frac{\lambda}{2}\sum_i\sum_k(x_k^{(i)})^2.
\label{eq:rec-cost}
\end{equation}

Unlike supervised learning, both user parameters $\{(\vec{w}^{(j)},b^{(j)})\}$
and item features $\{\vec{x}^{(i)}\}$ are simultaneously optimised.

\subsection{Mean Normalisation}

A new user with no ratings yields $\vec{w}^{(j)}=\mathbf{0}$,
$\hat{y}^{(i,j)}=0$ --- unhelpful. Mean normalisation:
\[\mu_i=\frac{\sum_{j:\,r^{(i,j)}=1}y^{(i,j)}}{\sum_j r^{(i,j)}},\qquad
y_{\text{norm}}^{(i,j)}=y^{(i,j)}-\mu_i.\]
Prediction: $\hat{y}^{(i,j)}=\vec{w}^{(j)}\cdot\vec{x}^{(i)}+b^{(j)}+\mu_i$.
A new user now gets $\hat{y}=\mu_i$ (average rating) --- a sensible default.

\subsection{Finding Related Items}

Learned item features $\vec{x}^{(i)}$ can be used to find similar items:
\[\mathrm{sim}(i,i')=\|\vec{x}^{(i)}-\vec{x}^{(i')}\|_2^2.\]

%----------------------------------------------------------
\section{Content-Based Filtering}
\label{sec:cbf}

\subsection{Overview}

Content-based filtering recommends items by matching item features to a model
of the user's preferences. It does not need other users' data.

\subsection{User Profile and Prediction}

Let $\vec{x}^{(i)}\in\R^d$ be the feature vector for item $i$. A linear user
profile $\vec{\theta}^{(j)}$ is estimated from rating history:
\[\hat{y}^{(i,j)}=\vec{\theta}^{(j)}\cdot\vec{x}^{(i)}+b^{(j)}.\]

\paragraph{Advantages.}
\begin{itemize}
\item Handles cold-start for new \emph{items} (as long as features are available).
\item Explainable: ``recommended because you liked action movies.''
\item No popularity bias.
\end{itemize}

\paragraph{Limitations.}
\begin{itemize}
\item Requires rich item features.
\item Over-specialisation: keeps recommending similar content.
\item User cold-start remains.
\end{itemize}

\subsection{Hybrid Systems}

Production systems typically combine collaborative and content-based filtering.
Google's YouTube recommender uses a two-tower neural network: one tower
processes user features (including watch history), the other processes video
content features; the final score is the dot product of the two embeddings.

%----------------------------------------------------------
\section{Principal Component Analysis}
\label{sec:pca}

\subsection{Motivation}

Real-world data often contain far fewer independent sources of variation than
the raw number of features suggests. PCA finds the directions of maximum
variance and projects data onto a lower-dimensional subspace.

\subsection{Mathematical Formulation}

Let $\mathbf{X}\in\R^{n\times d}$ be mean-centred. The covariance matrix:
\[\boldsymbol{\Sigma}=\frac{1}{n-1}\mathbf{X}^{\top}\mathbf{X}.\]

The $l$-th principal component is the eigenvector of $\boldsymbol{\Sigma}$
with the $l$-th largest eigenvalue $\lambda_l$. The projection:
\[\mathbf{Z}=\mathbf{X}\mathbf{U}_k,\]
where $\mathbf{U}_k\in\R^{d\times k}$ contains the top-$k$ eigenvectors.

Equivalently via SVD: $\mathbf{X}=\mathbf{U}\mathbf{S}\mathbf{V}^{\top}$;
principal components are the columns of $\mathbf{V}$.

\subsection{Choosing the Number of Components}

Fraction of variance explained:
\[\mathrm{EVR}(k)=\frac{\sum_{l=1}^k\lambda_l}{\sum_{l=1}^d\lambda_l}.\]
Common heuristic: choose $k$ such that $\mathrm{EVR}(k)\ge 0.95$.

\begin{figure}[h]
\centering
\begin{tikzpicture}
\begin{axis}[
  width=0.6\textwidth, height=5cm,
  xlabel={Principal component}, ylabel={Variance explained (\%)},
  xmin=0, xmax=8, ymin=0, ymax=58,
  xtick={1,2,3,4,5,6,7},
  grid=both, grid style={gray!15, thin},
  tick label style={font=\small}, label style={font=\small},
  axis lines=left,
  every axis plot/.append style={line width=1.4pt}]
\addplot[pBlue, mark=*, mark size=2.5pt]
  coordinates{(1,50)(2,20)(3,12)(4,7)(5,4)(6,3)(7,2)};
\addplot[pOrange, dashed, mark=*, mark size=2.5pt]
  coordinates{(1,50)(2,70)(3,82)(4,89)(5,93)(6,96)(7,98)};
\draw[pRed, dashed, thin] (axis cs:0,50) -- (axis cs:8,50);
\node[right, font=\small, pRed] at (axis cs:0.1,52)
  {PC1: 50\% of variance};
\end{axis}
\end{tikzpicture}
\caption{Scree plot (blue, left axis) and cumulative variance explained (orange,
right axis). Most variance is captured by the first few components.}
\label{fig:scree}
\end{figure}

\subsection{PCA in Recommender Systems}

\begin{enumerate}
\item \textbf{Dimensionality reduction}: compress high-dimensional content
  features before feeding into a model.
\item \textbf{Latent Semantic Analysis}: truncated SVD directly on the
  user--item matrix fills in missing ratings:
  $\hat{\mathbf{R}}=\mathbf{U}_k\mathbf{S}_k\mathbf{V}_k^{\top}$.
\item \textbf{Visualisation}: project learned embeddings to 2D/3D.
\item \textbf{Noise reduction}: discard small singular values.
\end{enumerate}

\paragraph{Connection to matrix factorisation.}
For a fully observed matrix with no regularisation, the rank-$k$ matrix
factorisation recovers exactly the truncated SVD:
\[\underset{\mathbf{U},\mathbf{V}}{\argmin}
\|\mathbf{R}-\mathbf{U}\mathbf{V}^{\top}\|_F^2
=\mathbf{U}_k\mathbf{S}_k\mathbf{V}_k^{\top}.\]

%----------------------------------------------------------
\section*{Chapter Summary}
\addcontentsline{toc}{section}{Chapter Summary}

\begin{enumerate}
\item \textbf{Collaborative filtering} exploits shared preferences. Matrix
  factorisation learns compact latent representations optimised by SGD or ALS.
\item \textbf{Mean normalisation} gives sensible default predictions for new
  users.
\item \textbf{Content-based filtering} builds a per-user preference model from
  item features; handles item cold-start but can over-specialise.
\item \textbf{PCA} extracts the principal axes of variation; it underpins matrix
  factorisation and is used for dimensionality reduction, noise filtering, and
  latent semantic analysis.
\item Production systems \textbf{combine} both paradigms in large-scale deep
  learning architectures.
\end{enumerate}

% =========================================================
%  Chapter 10: Reinforcement Learning
% =========================================================
\chapter{Reinforcement Learning}

Reinforcement Learning (RL) is the third major pillar of machine learning,
alongside supervised and unsupervised learning. Where supervised learning
requires labelled examples of the correct answer and unsupervised learning
discovers hidden structure, RL trains an \textbf{agent} to take sequential
decisions by interacting with an \textbf{environment} and receiving feedback
as numerical \textbf{rewards}. The key feature: the system is told
\emph{what to achieve} (via the reward function), not \emph{how to achieve it}.

%----------------------------------------------------------
\section{Core Concepts}

\subsection{What Is Reinforcement Learning?}

Learning proceeds by \textbf{trial and error}: the agent tries actions,
observes consequences, receives reward/penalty signals, and gradually improves
its behaviour to maximise cumulative reward.

\textit{Analogy.} Training a dog: sitting on command earns a treat (positive
reward); misbehaving earns a scolding (negative reward). Over many repetitions
the dog learns which behaviours earn treats.

\paragraph{RL vs.\ supervised learning.}
\begin{center}
\begin{tabular}{lll}
\toprule
\textbf{Feature} & \textbf{Supervised Learning} & \textbf{Reinforcement Learning}\\
\midrule
Input/output & $x\mapsto y$ (label) & $s\mapsto a$ (action)\\
Data & Expert-labelled dataset & Reward function + experience\\
Objective & Predict correct label & Maximise cumulative reward\\
Exploration & None & Essential\\
\bottomrule
\end{tabular}
\end{center}

A fundamental advantage of RL for control tasks: for complex motor control
(robotic locomotion, aerobatic flight) it is often impossible to specify the
optimal action in every state. A reward function sidesteps this.

\subsection{Key Components}

\begin{enumerate}
\item \textbf{State} $s$: complete description of the current situation.
\item \textbf{Action} $a$: the decision made at each step.
\item \textbf{Reward} $R(s)$: scalar feedback signal.
\item \textbf{Policy} $\pi$: function mapping states to actions,
  $\pi(s)=a$.
\item \textbf{Return} $G$: discounted cumulative reward (defined below).
\end{enumerate}

\subsection{The Discount Factor and Return}

The \textbf{discount factor} $\gamma\in[0,1)$ encodes the agent's degree of
patience.

\begin{equation}
\boxed{G=R_1+\gamma R_2+\gamma^2 R_3+\cdots
=\sum_{k=0}^\infty\gamma^k R_{k+1}.}
\label{eq:return}
\end{equation}

\begin{itemize}
\item $\gamma\approx 1$ (patient): future rewards nearly as valuable as
  immediate ones.
\item $\gamma\approx 0$ (myopic): rewards a few steps away are heavily
  discounted.
\end{itemize}

\textit{Financial analogy.} $\gamma$ mirrors the time value of money: a
reward now is more valuable than the same reward after $k$ time steps.

\paragraph{Example (Mars Rover, $\gamma=0.5$).}
Starting at state $s=4$, moving left: reward sequence $[0,0,0,100]$.
\[G=0+0.5(0)+0.25(0)+0.125(100)=12.5.\]

\begin{center}
\begin{tabular}{llcc}
\toprule
Start & Policy & Reward sequence & Return ($\gamma=0.5$)\\
\midrule
$s=4$ & Always Left  & $[0,0,0,100]$ & $12.5$\\
$s=2$ & Always Left  & $[0,100]$     & $50.0$\\
$s=4$ & Always Right & $[0,0,40]$    & $10.0$\\
$s=5$ & Always Right & $[0,40]$      & $20.0$\\
\bottomrule
\end{tabular}
\end{center}

\subsection{Markov Decision Processes (MDPs)}

The standard mathematical framework for RL.

\begin{itemize}
\item \textbf{Markov property}: the future depends only on the present state,
  not on the history.
\item \textbf{Agent--environment loop}: observe $s$, select $a=\pi(s)$,
  environment transitions to $s'$, agent receives $R(s)$. Repeat.
\end{itemize}

\begin{figure}[h]
\centering
\begin{tikzpicture}[
  box/.style={draw, rounded corners=4pt, minimum width=2.8cm,
              minimum height=1.0cm, align=center, fill=nodeGray!60},
  arr/.style={-{Stealth[length=3mm]}, thick}]
\node[box] (agent) {Agent\\$\pi(s)=a$};
\node[box, right=4.5cm of agent] (env) {Environment};
\draw[arr] (agent.east) -- node[above]{\small action $a$} (env.west);
\draw[arr] (env.south) -- ++(0,-0.9) -|
  node[below, pos=0.25]{\small reward $R(s)$} (agent.south);
\draw[arr] (env.north) -- ++(0,0.9) -|
  node[above, pos=0.25]{\small new state $s'$} (agent.north);
\end{tikzpicture}
\caption{The agent--environment interaction loop in a Markov Decision Process.}
\label{fig:mdp-loop}
\end{figure}

%----------------------------------------------------------
\section{State-Action Value Function and the Bellman Equation}

\subsection{The Q-Function}

\begin{definition}[Q-function / Action-value function]
$Q(s,a)$ is the \textbf{total discounted return} obtained by:
\begin{enumerate}
\item Starting in state $s$.
\item Taking action $a$ exactly once.
\item Following the \emph{optimal} policy thereafter.
\end{enumerate}
\end{definition}

The optimal policy selects:
\begin{equation}
\pi^*(s)=\argmax_a\,Q(s,a).
\label{eq:greedy}
\end{equation}

\begin{center}
\begin{tabular}{llcc}
\toprule
State $s$ & Action $a$ & Reward sequence & $Q(s,a)$ ($\gamma=0.5$)\\
\midrule
$s=2$ & Left  & $[0,100]$  & $0.5\times 100=50.0$\\
$s=2$ & Right & $[0,0,0,100]$ & $0.5^3\times 100=12.5$\\
$s=4$ & Left  & $[0,0,0,100]$ & $12.5$\\
$s=4$ & Right & $[0,0,40]$    & $10.0$\\
\bottomrule
\end{tabular}
\end{center}

\subsection{The Bellman Equation}

\begin{equation}
\boxed{Q(s,a)=R(s)+\gamma\max_{a'}Q(s',a').}
\label{eq:bellman}
\end{equation}

\paragraph{Components.}
\begin{description}
\item[$R(s)$] Immediate reward for the current step.
\item[$\gamma$] Discount factor.
\item[$\max_{a'}Q(s',a')$] Best possible future return from the next state.
\end{description}

At terminal states: $Q(s,a)=R(s)$.

\paragraph{Derivation sketch.}
\[G=R_1+\gamma R_2+\gamma^2 R_3+\cdots
=R(s)+\gamma(R_2+\gamma R_3+\cdots)
=R(s)+\gamma\max_{a'}Q(s',a').\quad\square\]

\paragraph{Stochastic environments.}
When transitions are random, the Bellman equation includes an expectation:
\begin{equation}
Q(s,a)=R(s)+\gamma\,\mathbb{E}[\max_{a'}Q(s',a')].
\label{eq:bellman-stoch}
\end{equation}

\begin{lstlisting}[caption={Tabular value iteration for the Mars Rover}]
import numpy as np

num_states, gamma = 6, 0.5
rewards   = np.array([100, 0, 0, 0, 0, 40])
q_table   = np.zeros((num_states, 2))   # 2 actions: left=0, right=1

def next_state(s, a):
    return max(0, s-1) if a == 0 else min(5, s+1)

for _ in range(20):
    for s in range(1, 5):
        for a in range(2):
            s2 = next_state(s, a)
            terminal = (s2 in [0, 5])
            r = rewards[s2] if terminal else 0
            q_table[s, a] = r + (0 if terminal else gamma*q_table[s2].max())

policy = {s: ['Left','Right'][q_table[s].argmax()] for s in range(1,5)}
print(policy)
\end{lstlisting}

%----------------------------------------------------------
\section{Continuous State Spaces and Deep Q-Networks}

\subsection{Continuous States}

Real problems require continuous state vectors. A self-driving truck:
$s=[x,y,\theta,\dot{x},\dot{y},\dot{\theta}]$. A helicopter: 12 state
variables. A lookup table is no longer feasible; a neural network must
approximate $Q(s,a)$.

\subsection{Lunar Lander Case Study}

\textbf{State}: $s=[x,y,\dot{x},\dot{y},\theta,\dot{\theta},l,r]$ (8D).

\textbf{Actions}: do nothing, fire left thruster, fire main engine, fire right
thruster.

\begin{center}
\begin{tabular}{lcc}
\toprule
\textbf{Event} & \textbf{Reward} & \textbf{Purpose}\\
\midrule
Reaching landing pad & $+100$ to $+140$ & Primary objective\\
Moving toward pad     & Positive          & Progress\\
Crash                & $-100$            & Severe penalty\\
Soft landing         & $+100$            & Bonus\\
Each leg grounded    & $+10$             & Stability\\
Main engine firing   & $-0.3$            & Fuel cost\\
\bottomrule
\end{tabular}
\end{center}

Learning objective: maximise $G_t=\sum_{k\ge 0}\gamma^k R_{t+k+1}$,
$\gamma=0.985$.

\subsection{DQN Architecture}

Feed only the state $s$ (8 values) into the network; output
\emph{all four Q-values simultaneously}:

\begin{itemize}
\item \textbf{Input}: 8D state vector.
\item \textbf{Hidden}: two fully connected layers, 64 ReLU units each.
\item \textbf{Output}: 4 scalar units, one per action.
\end{itemize}

\textbf{Advantage}: a single forward pass identifies $\argmax_a Q(s,a)$,
replacing four separate passes.

\paragraph{Training target.}
Treat as supervised learning: input $=(s,a)$, target $=R(s)+\gamma\max_{a'}Q(s',a')$.

\subsection{The DQN Algorithm}

\begin{enumerate}
\item Initialise network $Q$ with random weights.
\item Collect experience: store $(s,a,R(s),s')$ tuples in a
  \textbf{replay buffer}.
\item Sample a mini-batch from the buffer.
\item Compute targets $y_i=R(s_i)+\gamma\max_{a'}Q(s_i',a')$.
\item Train $Q_{\text{new}}$ on the mini-batch using MSE loss.
\item Soft-update: $Q\leftarrow Q_{\text{new}}$.
\item Repeat.
\end{enumerate}

\subsection{The $\varepsilon$-Greedy Policy}

With probability $\varepsilon$: take a random action (exploration).
With probability $1-\varepsilon$: take $\argmax_a Q(s,a)$ (exploitation).

\textbf{Decay schedule}: start $\varepsilon=1.0$, decay toward $0.01$ over
training.

\begin{figure}[h]
\centering
\begin{tikzpicture}
\begin{axis}[
  width=10cm, height=5.5cm,
  xlabel={Training episode}, ylabel={$\varepsilon$},
  xmin=0, xmax=1000, ymin=0, ymax=1.1,
  grid=major, grid style={gray!15, thin},
  tick label style={font=\small}, label style={font=\small},
  every axis plot/.append style={line width=1.4pt}]
\addplot[pBlue, domain=0:1000, samples=200]
  {0.99*exp(-0.005*x)+0.01};
\addplot[pGray, dashed] coordinates{(0,0.01)(1000,0.01)};
\node[font=\small, pGray] at (axis cs:600,0.07)
  {$\varepsilon_{\min}=0.01$};
\end{axis}
\end{tikzpicture}
\caption{Exponential $\varepsilon$-decay from 1.0 toward 0.01 over training.}
\label{fig:eps-decay}
\end{figure}

\subsection{Mini-Batch and Soft Updates}

\paragraph{Mini-batch gradient descent.}
Sample $m'\ll m$ examples from the replay buffer per update step. Cheaper per
step; more steps per wall-clock time; noisy but converges faster in practice.

\paragraph{Soft updates.}
Blend new parameters gradually:
\[W\leftarrow\tau W_{\text{new}}+(1-\tau)W,\quad
  b\leftarrow\tau b_{\text{new}}+(1-\tau)b.\]
Typical $\tau=0.01$: only $1\%$ of new parameters incorporated per step.
Prevents a single poor mini-batch from corrupting the policy.

\begin{center}
\begin{tabular}{lll}
\toprule
\textbf{Technique} & \textbf{Purpose} & \textbf{Benefit}\\
\midrule
Replay buffer    & Decouple collection from training & Data diversity\\
Mini-batching    & Updates on random subsets & Training speed\\
Soft updates     & Blend new weights gradually & Stability\\
$\varepsilon$-greedy & Balance exploration/exploitation & Avoid local optima\\
\bottomrule
\end{tabular}
\end{center}

%----------------------------------------------------------
\section{The State of Reinforcement Learning}

\paragraph{Current deployment.}
Despite breakthroughs (AlphaGo, AlphaStar, OpenAI Five), RL deployment in
industry is narrower than supervised/unsupervised learning. Most production ML
uses labelled data.

\paragraph{The simulation-to-reality gap.}
RL algorithms are typically developed in simulation. Policies trained there may
rely on unrealistic assumptions and fail on real hardware (e.g.\ sensor noise,
actuator delays). Active research in domain randomisation and sim-to-real
transfer aims to close this gap.

\paragraph{Future directions.}
\begin{itemize}
\item \textbf{Autonomous control}: the primary framework for sequential
  decision-making in robotics, logistics, and finance.
\item \textbf{Alignment research}: RLHF (Reinforcement Learning from Human
  Feedback) is central to training large language models to follow human intent.
\end{itemize}

%----------------------------------------------------------
\section*{Chapter Summary}
\addcontentsline{toc}{section}{Chapter Summary}

\begin{enumerate}
\item RL trains an agent to maximise cumulative discounted reward via trial and
  error, without requiring labelled examples of correct actions.
\item Key components: state $s$, action $a$, reward $R(s)$, discount $\gamma$,
  policy $\pi$, return $G=\sum\gamma^k R_{k+1}$.
\item An MDP formalises the problem under the Markov property.
\item The Q-function $Q(s,a)$ gives the optimal return for taking action $a$
  in state $s$ then acting optimally. $\pi^*(s)=\argmax_a Q(s,a)$.
\item The \textbf{Bellman equation} $Q(s,a)=R(s)+\gamma\max_{a'}Q(s',a')$
  is the recursive foundation of all Q-learning algorithms.
\item In stochastic environments: $Q(s,a)=R(s)+\gamma\mathbb{E}[\max_{a'}Q(s',a')]$.
\item Deep Q-Networks use a neural network to represent $Q$ over continuous
  state spaces.
\item Key engineering: replay buffer, mini-batch updates, soft target updates,
  $\varepsilon$-greedy exploration with decay.
\item RL deployment remains limited by the simulation-to-reality gap, but
  RLHF is now critical for aligning large language models.
\end{enumerate}


\end{document}
